\documentclass[11pt]{amsart}
\usepackage{amsmath,amssymb,amsthm}
\usepackage{booktabs}
\usepackage[margin=1.15in]{geometry}
\usepackage[colorlinks=true,linkcolor=blue,citecolor=blue,urlcolor=blue]{hyperref}
\usepackage{microtype}
\emergencystretch=1.5em

\newtheorem{theorem}{Theorem}[section]
\newtheorem{proposition}[theorem]{Proposition}
\newtheorem{corollary}[theorem]{Corollary}
\newtheorem{lemma}[theorem]{Lemma}
\newtheorem{conjecture}[theorem]{Conjecture}
\newtheorem{problem}[theorem]{Problem}
\newtheorem{question}[theorem]{Question}
\theoremstyle{definition}
\newtheorem{definition}[theorem]{Definition}
\newtheorem{example}[theorem]{Example}
\theoremstyle{remark}
\newtheorem{remark}[theorem]{Remark}
\newtheorem{observation}[theorem]{Observation}

\newcommand{\PP}{\mathbb{P}}
\newcommand{\CC}{\mathbb{C}}
\newcommand{\ZZ}{\mathbb{Z}}
\newcommand{\QQ}{\mathbb{Q}}
\newcommand{\RR}{\mathbb{R}}
\newcommand{\NN}{\mathbb{N}}
\newcommand{\gfrak}{\mathfrak{g}}
\newcommand{\Dscr}{\mathfrak{D}}
\newcommand{\chat}{\hat{c}}
\newcommand{\Bal}{\operatorname{Balanced}}
\newcommand{\Crit}{\operatorname{Crit}}
\newcommand{\Mbar}{\overline{\mathcal{M}}}
\newcommand{\Ascr}{\mathcal{A}}
\newcommand{\Om}{\Omega_{r,s}}

\hyphenation{scat-ter-ing tropi-cal tropi-cal-izes singu-lar-ity de-scen-dent}

\title[The scattering structure of open FJRW theory of $x^r+y^s$]{The scattering
structure of open FJRW theory of $x^{r}+y^{s}$}

\author{Bernd Johannes Wuebben}
\subjclass[2020]{14N35 (primary); 14J33, 14T90, 53D45 (secondary)}
\keywords{Open FJRW theory, scattering diagrams, wall-crossing, tropical
vertex, Landau--Ginzburg models, broken lines}
\date{August 15, 2026}

\begin{document}

\begin{abstract}
Gross, Kelly and Tessler constructed a genus-zero open enumerative theory for
the Landau--Ginzburg model $(x^{r}+y^{s},\mu_{r}\times\mu_{s})$: the open FJRW
invariants depend on boundary conditions, and the systems of invariants form a
torsor under a wall-crossing group which they observe is ``very similar in
spirit'' to the tropical vertex group of Kontsevich--Soibelman. We turn the
observation into a structure theory. Every primary transition between invariant
systems has a unique slope-ordered factorization by the
boundary rays, to every order in the deformation parameters; for descendents the
compatible factors are defined coefficientwise
on finite divisor-closed coefficient sets. They define a relative scattering structure on an
integral-affine sector of directions. Its pro-nilpotent walls are supported precisely on the deformation
monomials that carry a critical graph of positive Euler weight, while
marginal critical graphs supply a torus factor,
sharpening the central-charge threshold for wall-crossing; two unique
extremal systems of invariants produce a canonical relative scattering
diagram, and the open
topological recursion supplies the direction transverse to its walls, so that
the closed theory determines all invariants up to the torsor action. Once
descendents are admitted the walls accumulate densely, so open chambers are
replaced by coefficientwise slope cuts. The potential
family is then the coefficientwise broken-line transport of a unique cut-independent seed
system, for a canonical assignment of each potential coefficient to a slope cut. The seeds
admit a finite ordered-chain inversion formula and are virtual Euler degrees of compact
bar-cascade objects assembled from relative-Euler and critical-homotopy zero loci:
the seed values vanish at the extreme coefficients that receive no bending
contributions from other initial monomials, a consequence delimited by the
interlacing of walls and coefficients and by a monotonicity property of bending
(the purely tropical part of the theory, where the
invariants are genuine broken-line counts), while the nonzero seed values are
non-tropical input, beyond what the diagram determines.
Every finite divisor-closed primary or descendent coefficient window is realized,
in increasing slope, by a concatenated family of simple canonical homotopies of
boundary conditions. Relative versions of the underlying extension arguments
make these realizations compatible coefficientwise; staging is forced by a
provable obstruction. The two extreme indices admit symmetric polar
representatives: on every wall-bearing diagonal, all balanced sections except
the surviving endpoint section are nowhere zero.  The resulting
homotopy-coherent transport on the punctured boundary-charge sector is
projective: it factors through the interval of slopes, so radial paths act
trivially.  Since the punctured sector is contractible, its total
extreme-to-extreme product is boundary-to-boundary transport rather than
monodromy about the common ray origin.  This gives a precise obstruction to
recovering a genuine singular affine position space from the present
continuation data without an additional valuation or affine gluing.
For $x^{5}+y^{5}$ we compute the first wall
functions and several new invariants, and verify the recursion in full against
exactly computed closed data. Three corrections to the literature accompany the
results. All computations are exact, and code reproducing them accompanies the
paper.
\end{abstract}

\maketitle
\enlargethispage{3pt}

\section{Introduction}
\label{sec:intro}

\subsection{The theory}
Gross, Kelly and Tessler \cite{GKT} defined genus-zero open FJRW invariants
\[
\nu_{\Gamma,\mathbf d}
=\Big\langle\prod_{i}\tau_{d_{i}}^{(a_{i},b_{i})}\,
\sigma_{1}^{k_{1}}\sigma_{2}^{k_{2}}\sigma_{12}\Big\rangle^{\mathfrak s,o}
\]
for the Landau--Ginzburg model $(W=x^{r}+y^{s},\ \mu_{r}\times\mu_{s})$: relative
Euler numbers of descendent Witten bundles over moduli of $W$-spin disks, which are
real orbifolds with corners. The invariants depend on a family $\mathfrak s$
of canonical multisections (boundary conditions). Three of their results organize this
dependence completely. First, the systems of invariants arising from canonical
families form a set $\Om(I)$ ($I$ a finite set of internal marking labels as
in \cite{GKT}; in universal finite-order statements below it is chosen
\emph{saturated} at the order under consideration, meaning it carries at least
$k-1$ labels of every primary twist at order $k$), characterized by explicit polynomial
conditions (the \emph{chamber-index} conditions of \cite[Def.~3.46,
Cor.~5.13]{GKT}; a \emph{chamber index} is their name for such a system of
invariants, not a chamber of any diagram in this paper):
normalizations for at most one internal marking, and the conditions
$\Ascr(J,\mathbf d,\nu)=(-1)^{d}$ (singletons) and $\Ascr(J,\mathbf d,\nu)=0$ for
$d(J,\mathbf d)<0$ (with $d(J,\mathbf d)$ as in \cite[Def.~4.8]{GKT}), where
$\Ascr$ is a $\Gamma$-function-weighted polynomial in the invariants. Second, $\Om(I)$ is a torsor: a wall-crossing group acts on it
faithfully and transitively \cite[Thm.~4.26]{GKT}, with one generating summand per
\emph{critical graph}, the disk type responsible for jumps. Third, the
generating-function potential $W^{\nu}$ of any system has oscillatory integrals
over the good basis equal to the closed extended FJRW generating function
\cite[Thm.~0.3]{GKT}, an open mirror theorem building on the closed
Landau--Ginzburg mirror symmetry of \cite{LLSS,HLSW}.

Gross--Kelly--Tessler remark \cite[Rem.~0.8]{GKT} that their wall-crossing group is
close in spirit to the tropical vertex group of Kontsevich--Soibelman \cite{KS06},
as used for Fano mirror symmetry in \cite{GPS10,Gross10} and in the log
Calabi--Yau setting \cite{GHK15}, and their survey raises
the problem of a tropical formulation of the open theory \cite{GKTsurvey}. Such a
formulation is not immediate. A Landau--Ginzburg model carries no target space and
hence no SYZ fibration, and on these grounds Maher \cite[Rem.~5.1.1]{Maher} records
the combinatorial machinery of the Gross--Siebert program as out of reach for these
models. The present paper develops the scattering structure of the theory and
computes its first walls. The construction is algebraic, carried out over the
wall-crossing group rather than over a degenerating target, which is what makes it
available here. Every finite primary or descendent coefficient window is
realized by canonical homotopies in Theorem~\ref{thm:georealization}, compatibly
in the coefficientwise limit.  The sector itself is the effective cone of the
two open-FJRW boundary charges.  Theorem~\ref{thm:positionobstruction} proves
that the induced continuation theory on its punctured cone factors through
slope and has no origin loop.  Thus the cone is an intrinsic relative
scattering sector, but a singular affine position space of the kind supplied by
an SYZ fibration requires genuinely additional input.

\subsection{Results}
Throughout, $\gfrak^{r,s}$ denotes the Landau--Ginzburg wall-crossing Lie algebra
of \cite[Def.~4.22]{GKT}: it is spanned by elements $u_{J,\mathbf d}X_{k_{1},k_{2}}$
with
\[
X_{k_{1},k_{2}}=x^{k_{1}}y^{k_{2}}\bigl((k_{2}{+}1)x\partial_{x}-(k_{1}{+}1)
y\partial_{y}\bigr)
\]
the Hamiltonian vector field of $x^{k_{1}+1}y^{k_{2}+1}$ for $dx\wedge dy$, and
$u_{J,\mathbf d}=\prod_{j}t_{a_{j},b_{j},d_{j}}$ a monomial in the deformation
parameters ($t_{a,b,d}$ dual to the insertion $\tau_{d}^{(a,b)}$; the
descendent index is suppressed when $d=0$), subject to the
congruences $k_{1}\equiv r(J)\ (\mathrm{mod}\ r)$, $k_{2}\equiv s(J)\
(\mathrm{mod}\ s)$ and the dimension constraint
\begin{equation}
\label{eq:dimconstraint}
sk_{1}+rk_{2}=m(J,\mathbf d)-rs,
\qquad
m(J,\mathbf d)=rs+\sum_{j}\bigl(sa_{j}+rb_{j}+rs(d_{j}-1)\bigr).
\end{equation}
The right-hand side of \eqref{eq:dimconstraint} is the \emph{Euler weight}
$w(J,\mathbf d)$ of the monomial $u_{J,\mathbf d}$. By the classification
\cite[Not.~2.38]{GKT}, each $(J,\mathbf d)$ carries an integer count
$N=N(J,\mathbf d)$ (Definition~\ref{def:diagonal}): when $N\ge0$ there are $N{+}1$ balanced disk types
(potential coefficients), and when $N\ge1$ there are $N$ critical graphs
(walls); below the respective thresholds those sets are empty. The Lie
bidegree of a critical graph is its boundary degree shifted by $(-1,-1)$: the wall
Hamiltonian $x^{k_{1}+1}y^{k_{2}+1}$ is the boundary monomial
$\sigma_{1}^{\#}\sigma_{2}^{\#}$ of the critical disk, and the shift is the removal
of the single fully twisted point $\sigma_{12}$ that balanced disks carry and
critical disks do not. This is the primitive form $dx\wedge dy$ acting as an
integral-affine shift (\S\ref{sec:algebra}).

\begin{theorem}[Weight trichotomy and the threshold; \S\ref{sec:algebra}]
\label{thm:trichotomy}
In the enlarged algebra of Remark~\ref{rem:torus} (the wall-crossing algebra
with the marginal torus directions adjoined), for a nonempty marking
multiset $J$, wall summands over the monomial
$u_{J,\mathbf d}$ exist if and only if $N(J,\mathbf d)\ge1$ and
$w(J,\mathbf d)\ge0$. The pro-nilpotent summands are exactly those of positive
weight. For a pro-nilpotent wall, positive weight is necessary but not sufficient (at $r=s=5$ the
monomial $t_{3,3}$ has $w=5$ but $N=0$, hence no wall); at $w=0$ the constraint
\eqref{eq:dimconstraint} admits only the marginal torus bidegree $(0,0)$, and at
$w<0$ it admits nothing. Consequently a monomial of nonpositive weight carries
no pro-nilpotent wall generator of its own (although higher-order transport
from divisor monomials may reach its coefficients), and the central charge
$\chat_{W}=2-\tfrac2r-\tfrac2s$
governs a trichotomy: pro-nilpotent $d=0$ wall-crossing exists if and only if
$\chat_{W}>1$ (we call the pair $(r,s)$ \emph{hyperbolic}); at $\chat_{W}=1$
(\emph{parabolic}) primary wall-crossing exists but is carried entirely by the
marginal torus direction; for $\chat_{W}<1$ (\emph{simple}) there is no
primary wall-crossing at all.
\end{theorem}

\begin{theorem}[The relative scattering presentation; \S\ref{sec:presentation}]
\label{thm:presentation}
(i) (Loop consistency.) For any family of simple
canonical homotopies supplied by \cite[Lem.~3.39]{GKT} with equal endpoints,
the time-ordered
product of the critical-graph jump automorphisms constructed in the proof of
\cite[Thm.~5.8]{GKT} is the
identity; for such homotopies with arbitrary endpoints, the total product depends
only on the endpoints. (ii) (Boundary-ray presentation.) In the primary
algebra, modulo every power of the deformation ideal, every element of the
enlarged wall-crossing group factors uniquely as a slope-ordered product of
exponentials supported on the shifted boundary rays
$\RR_{\ge0}(k_{1}{+}1,k_{2}{+}1)$; the marginal torus terms belong to the
$(1,1)$-ray. The compatible factorizations define its completed boundary-ray
presentation. In the full descendent algebra the same presentation is taken on
finite divisor-closed coefficient sets and passed to the coefficientwise limit.
We call the resulting ray-decorated,
boundary-to-boundary transport a \emph{relative scattering structure}: unlike a
standard consistent diagram on a complete affine base, its total transport need
not be the identity.
\end{theorem}

The statements below use the following language. The \emph{diagonal} of a
deformation monomial is its set of balanced coefficients, the lattice points
of one line in the boundary-degree plane (Definition~\ref{def:diagonal}).
The rest is developed in
\S\ref{sec:geometric}: the \emph{direction sector} is the closed first
quadrant of boundary-monomial directions; a \emph{slope cut} is a direction in
the sector together with a choice of side when a wall sits on that direction;
the transport $\Psi(q)$ to a cut $q$ is the slope-ordered product of the wall
factors from the $\sigma_{1}$-axis side (the \emph{base cut}) up to $q$; and a
bending contribution to a coefficient is \emph{foreign} if it originates from
another monomial's seed (Definitions~\ref{def:sector} and~\ref{def:anchor}).

\begin{theorem}[Normal forms and the canonical diagram; \S\ref{sec:canonical}]
\label{thm:canonical}
On every wall-bearing diagonal, including a marginal diagonal, the wall action
on the $N{+}1$ potential coefficients is triangular; consequently there exist
unique chamber indices $\nu^{\min}$ and $\nu^{\max}$ in which, on every such
diagonal, only the extreme coefficient survives. The boundary-slope-ordered
factorization of the unique group element
carrying $W^{\nu^{\min}}$ to $W^{\nu^{\max}}$ is a canonical, choice-free
completed primary relative diagram $\Dscr^{r,s}_{0}$, computable order by order;
its descendent extension $\Dscr^{r,s}$ is interpreted coefficientwise on slope
cuts. We write $\theta_{\rho}$ for the factor of the diagram supported on the
ray $\rho$. For $x^{5}+y^{5}$
its walls begin
\begin{gather*}
\theta_{(1,2)}=\exp\bigl(\tfrac1{25}\,t_{2,3}t_{3,3}X_{0,1}+\cdots\bigr),\qquad
\theta_{(2,1)}=\exp\bigl(\tfrac1{25}\,t_{3,2}t_{3,3}X_{1,0}+\cdots\bigr),\\
\theta_{(1,1)}=\exp\bigl(\tfrac1{50}\,t_{3,3}^{2}X_{1,1}+\cdots\bigr).
\end{gather*}
The subscripts label primitive boundary directions. A wall factor collects all
critical-graph terms for which $(k_{1}{+}1,k_{2}{+}1)$ lies on that ray
(Proposition~\ref{prop:shift}).
\end{theorem}

\begin{theorem}[Effective determination; \S\ref{sec:diagonal}]
\label{thm:effective}
On each $(J,\mathbf d)$-diagonal the wall directions span, in the NS-visible
case (both residues in the Neveu--Schwarz range: $r(J)\le r-2$ and
$s(J)\le s-2$), exactly the kernel of the good-basis period functional, and the chamber-index
conditions of \cite{GKT} amount to one $\Gamma$/period-weighted linear condition
per diagonal with right-hand side determined by strictly lower order. Together with
the characterization $\Om(I)=\Upsilon_{r,s}(I)$ of \cite[Cor.~5.13]{GKT}
($\Upsilon_{r,s}(I)$ the set of solutions of the chamber-index conditions) this
determines all open FJRW invariants of $x^{r}+y^{s}$, up to the torsor action,
from the normalizations and the closed extended theory: the topological recursion
of \cite[Thm.~0.5]{GKT} appears as the direction transverse to the walls.
\end{theorem}

Theorem~\ref{thm:effective} is the algebraic form of the statement that
\emph{the scattering structure computes the open theory}. The worked case
$x^{5}+y^{5}$ (\S\ref{sec:canonical}, \S\ref{sec:x5y5}) produces, beyond the wall
functions above, the invariant relations and values
\begin{gather}
\langle\tau_{0}^{(3,3)}\tau_{0}^{(3,3)}\sigma_{1}\sigma_{2}^{6}\sigma_{12}\rangle
+\langle\tau_{0}^{(3,3)}\tau_{0}^{(3,3)}\sigma_{1}^{6}\sigma_{2}\sigma_{12}\rangle
=-\tfrac25,
\label{eq:firstwallrel}\\
\langle\tau_{0}^{(3,3)}\tau_{0}^{(3,3)}\tau_{0}^{(3,3)}
\sigma_{1}^{4}\sigma_{2}^{4}\sigma_{12}\rangle=\tfrac15,\qquad
\langle\tau_{0}^{(1,2)}\tau_{0}^{(2,1)}\tau_{0}^{(2,2)}\sigma_{12}\rangle
=-\tfrac1{25},
\label{eq:newinvariants}
\end{gather}
the first a wall relation (only the sum is invariant; the difference jumps across
the first symmetric pro-nilpotent primary wall $t_{3,3}^{2}X_{1,1}$), the
second and third absolute: they live on diagonals with no walls, the second on a
diagonal invisible to the good-basis periods. The seed relation (for the seed system of Theorem~\ref{thm:transport} below)
takes the uniform form
\begin{equation}
\label{eq:seed}
\frac{\nu_{(r,0)}}{r}+\frac{\nu_{(0,s)}}{s}=-\frac1{rs}
\qquad\text{(e.g. }\langle\sigma_{1}^{r}\sigma_{12}\cdots\rangle
+\langle\sigma_{2}^{r}\sigma_{12}\cdots\rangle=-\tfrac1r\ \text{at }r=s),
\end{equation}
for every two-insertion seed diagonal $\{(a_{1},b_{1}),(a_{2},b_{2})\}$,
$|J|=2$, with $a_{1}+a_{2}=r$, $b_{1}+b_{2}=s$ (Remark~\ref{rem:ex55}). Here
$\nu_{(k_{1},k_{2})}$ denotes the invariant with boundary profile
$\sigma_{1}^{k_{1}}\sigma_{2}^{k_{2}}\sigma_{12}$, and the suppressed
insertions are the two interior markings of the seed.

The last three results are geometric: the transport representation, the
finite-support and coefficientwise realization of the diagram by homotopies of
boundary conditions, and the exact position-space obstruction carried by that
realization.

\begin{theorem}[The transport representation; \S\ref{sec:geometric}]
\label{thm:transport}
Fix the \emph{anchor assignment}: each monomial
$u_{J,\mathbf d}\,x^{k_{1}}y^{k_{2}}$ of the potential is anchored at the
one-sided slope cut through its direction $(k_{1},k_{2})$, on the side toward
the $(1,1)$-direction. A constant monomial is anchored at the base cut; on
slope $1$ either one-sided convention gives the same coefficient. Then, by
induction in the divisibility order of deformation monomials: on every monomial
$u_{J,\mathbf d}$ the difference between the potential's coefficient vector and
the transported contributions of the seeds of the proper divisor monomials is
\emph{independent of the cut}, and defining the \emph{seed system} by these
constants, every cut potential is the coefficientwise transport of the seeds,
\[
W(q)=\sum_{m}\Psi(q)\,\Psi(q_{m})^{-1}(\mathrm{seed}_{m}),
\]
the sum over the monomials $m$ of the potential and $q_{m}$ the anchor cut
of $m$. On the coefficient of a fixed deformation monomial $u$, project every
factor to the finite divisor-closed quotient consisting of the divisors of $u$;
only finitely many projected boundary-ray factors are nontrivial, and their
slope-ordered product defines $\Psi(q)$ coefficientwise from the base cut to
$q$. The seed system
is the unique one with this property. On a coefficient receiving no foreign
bending, the seed equals its anchored value; for a foreign-free extreme
coefficient of a
wall-bearing diagonal that value is its vanishing extreme-index value
(Proposition~\ref{prop:interlace}), so every such coefficient
carries seed zero. Those
zero-seed coefficients form the purely tropical part of the theory, while the
nonzero seed values are non-tropical input.
\end{theorem}

Theorem~\ref{thm:transport} is the structural core of the tropical description
of the open theory: the invariants are computed by broken lines (monotone
transports bending at the walls of $\Dscr^{r,s}$) from canonical seed data.
The seeds play the role of the initial monomials of broken lines in
\cite{Gross10,CPS}; the word carries no cluster-theoretic meaning here.  The
ordered-chain formula and corrected relative-Euler cascade cycles below give
them a zero-dimensional enumerative interpretation.  The two extreme indices
have symmetric polar representatives whose nonextreme balanced sections are
nowhere zero (Theorem~\ref{prop:polar}), and the cascade cycles are the virtual
Euler cycles of single compact Kuranishi bar realizations
(Theorem~\ref{thm:virtualcascade}).  Theorem~\ref{thm:positionobstruction}
determines the exact boundary of the position-space reading: the geometric
continuation exists on the punctured sector but factors through slope and
cannot supply origin monodromy.  The whole is the
Landau--Ginzburg analogue of the tropical disk-counting picture of
\cite{Gross10,CPS}.

\begin{theorem}[Finite-support and pro-geometric realization; \S\ref{ssec:realization}]
\label{thm:georealization}
Let $S$ be a finite divisor-closed set of deformation monomials containing
$1$.  Choose a finite labeled set $I$ with at least
\[
n_{a,b}(S):=\max_{u\in S}\sum_{d\ge0}\operatorname{ord}_{t_{a,b,d}}(u)
\]
labels of every twist $(a,b)$ occurring in $S$.  The quotient diagram has
finitely many rays $\rho_1,\dots,\rho_n$, listed in increasing slope.  There is a
canonical homotopy of boundary conditions, obtained by
concatenating families of simple canonical homotopies in the sense of
\cite[\S3.4]{GKT}, from a family realizing $\nu_I^{\min}$ to one realizing
$\nu_I^{\max}$, the restrictions of the universal extreme indices, such that
after reduction to the quotient supported on $S$ every nonzero
jump is supported on a single boundary ray, the rays carrying
nonzero reduced jumps occur in increasing slope as the homotopy time
increases, and the total reduced jump on $\rho_{i}$ is the wall automorphism
$\theta_{\rho_{i}}$ of Theorem~\ref{thm:canonical}.  These realizations can be
chosen compatibly over a cofinal construction-closed exhaustion of the labeled
graph/descendent data.  The resulting coefficientwise homotopy has
pro-symmetric endpoints $\nu^{\min},\nu^{\max}$, and every fixed graph and
descendent vector sees only finitely many nonconstant ray stages.
\end{theorem}

The staging is forced: above the central-charge threshold a single normal-form
family of \cite[Lem.~3.39]{GKT} provably cannot separate the walls of a
diagonal carrying two of them, since they all jump at one time
(Proposition~\ref{prop:obstruction}), and the exact first order at which two
walls share a diagonal is computed for every hyperbolic pair
(Proposition~\ref{prop:parabolicgeom}, Remark~\ref{rem:thresholdgeom}).
Realizing each boundary-ray factor in its own stage is what evades the
obstruction.

\begin{theorem}[Projective coherence and the position-space obstruction;
  \S\ref{ssec:missing}]
\label{thm:positionobstruction}
For every finite divisor-closed coefficient set $S$, the staged realization
extends to a homotopy-coherent transport pseudofunctor on the punctured sector
$C_\partial^\times=C_\partial\setminus\{0\}$: points give canonical
multisection families, tame paths give concatenated canonical homotopies, and
tame disks give relative two-parameter families.  An oriented crossing of a
wall $\rho$ induces $\theta_{\rho,S}^{\pm1}$.  This pseudofunctor factors,
up to coherent equivalence, through
\[
 \pi:C_\partial^\times\longrightarrow\mathbb P(C_\partial),
\]
and the factorizations are compatible coefficientwise as $S$ grows.
Consequently every radial path acts trivially and dilation changes no GKT
transport observable.  Moreover
$C_\partial^\times\cong(0,\infty)\times[0,1]$ is contractible: the product
$\prod_\rho^\rightarrow\theta_\rho=g^{\min\to\max}$ is transport between the
two boundary sides, not origin monodromy.  Hence the present GKT continuation
data do not intrinsically determine the radially detected modulus, singular
affine extension, or vertex monodromy required of the proposed position space.
An affirmative Gross--Siebert/SYZ and theta-function theory requires additional
valuation, affine-gluing, and binary enumerative input.
\end{theorem}

\subsection{Three corrections to the literature}
Our computations bear on three points in \cite{GKT,GKTsurvey}, recorded here and in
Remarks~\ref{rem:ex55} and~\ref{rem:torus} with the supporting derivations.
First, imposing the open mirror theorem \cite[Thm.~0.3]{GKT} directly on the
oscillatory expansion of the potential --- the mechanism of their Theorem~0.2 ---
gives the seed values \eqref{eq:seed}; at $r=s=4$ this yields $-\tfrac14$ for the
sum in \cite[Ex.~5.5]{GKTsurvey}, where the value $-\tfrac18$ is stated. The
partition count in the expansion of $\Ascr$ forced by \cite[Thm.~4.9]{GKT} (ordered
partitions of the labeled marking set, with $1/h!$) also reproduces, through the
proven \cite[Thm.~5.4]{GKT}, the closed value
$\Ascr(\{(1,1),(1,1)\},0)=\langle\tau_{0}^{(1,1)}{}^{3}\rangle=1$ at $r=s=5$,
pinning the same count independently. Second, at wall-bearing marginal monomials
the critical graphs have torus bidegree $(0,0)$, which the pro-nilpotent algebra of
\cite[Def.~4.22]{GKT} excludes; the proofs of \cite[Thm.~5.8, Cor.~5.13]{GKT}
nevertheless use exactly these summands (through their $Q_{0}$-sets and
Lemma~5.16), so the group there should be read as enlarged by the marginal torus
summands, the elliptic wall-crossing group appearing in Maher's thesis
\cite[\S5.2.2]{Maher}. All statements in this paper use the enlarged group.
Third, in the same proof, equation (5.51) of \cite{GKT} carries the coefficient
$s\,k_{2}(\Lambda_{J,1})$ where the $p=0$ instance of its own (5.52) gives
$s\,k_{1}(\Lambda_{J,1})$; the latter is the one consistent with the action
\eqref{eq:wallvector} of the wall on $W_{0}$, whose $y^{s}$-term carries
$s(k_{1}{+}1)$, and the displayed choice of $h_{1,\mathbf d}$ inherits the same
slip. Only that display is affected: the inductive choice for $p\ge1$ uses the
correct form, and we use it throughout.

\subsection{Method and verification}
All algebraic and numerical claims are verified by exact polynomial and rational
arithmetic in the accompanying programs (\S\ref{sec:cert}), which are available,
together with this paper, at the public repository
\url{https://github.com/bwuebben/open-fjrw-scattering}. An expository
companion note on the central-charge threshold and the first-wall census is
\cite{note}; the closed FJRW input
for $x^{5}+y^{5}$ is computed by three independent, mutually consistent routes, including a
confirmation of every value through Witten's $r$-spin class in
\texttt{admcycles}, and the open topological recursion of \cite[Thm.~0.5]{GKT} is
verified in full against this data (\S\ref{sec:x5y5}).

\section{The wall-crossing algebra of critical graphs}
\label{sec:algebra}

We recall the objects of \cite[\S2.4, \S4.3]{GKT} in the notation used throughout.
Internal markings carry twists $(a,b)$, $0\le a\le r-2$, $0\le b\le s-2$, and
descendent exponents $d\ge0$; boundary markings are of three types,
$\sigma_{1}$ (twist $(r{-}2,0)$), $\sigma_{2}$ (twist $(0,s{-}2)$), and the fully
twisted $\sigma_{12}$ (twist $(r{-}2,s{-}2)$). For a multiset $J$ of internal
twists with descendents $\mathbf d$, set $r(J)\equiv\sum_{j}a_{j}$
$(\mathrm{mod}\ r)$, $s(J)\equiv\sum_{j}b_{j}$ $(\mathrm{mod}\ s)$, and
$m(J,\mathbf d)$ as in \eqref{eq:dimconstraint}.

\begin{definition}[{\cite[Not.~2.38]{GKT}}]
\label{def:diagonal}
Write $\sum_{j}a_{j}=r(J)+\ell_{1}r$, $\sum_{j}b_{j}=s(J)+\ell_{2}s$ and
$N=N(J,\mathbf d)=\ell_{1}+\ell_{2}-|J|+1+\sum_{j}d_{j}$. The \emph{balanced}
disk types over $(J,\mathbf d)$ are $\Gamma_{J,p}$, $0\le p\le N$, with boundary
degrees $(k_{1},k_{2})=(r(J)+pr,\ s(J)+(N{-}p)s)$ and one $\sigma_{12}$-point; the
\emph{critical} graphs are $\Lambda_{J,p}$, $1\le p\le N$, with boundary degrees
$(r(J)+(p{-}1)r+1,\ s(J)+(N{-}p)s+1)$ and no $\sigma_{12}$-point. We refer to the
set of balanced coefficients over a fixed $(J,\mathbf d)$ as its \emph{diagonal}.
The balanced range is empty for $N<0$, and the critical range is empty for
$N<1$. By \cite[Prop.~2.30]{GKT}, the diagonal is the set of lattice points of
$sk_{1}+rk_{2}=m(J,\mathbf d)$ obeying the congruences.
\end{definition}

\begin{proposition}[The primitive-form shift]
\label{prop:shift}
The Lie bidegree of the critical graph $\Lambda_{J,p}$, in the sense of
\eqref{eq:dimconstraint}, is its boundary degree minus $(1,1)$
(\cite[Prop.~2.37]{GKT}); equivalently, the wall element attached to $\Lambda$ is
the Hamiltonian vector field, for the primitive form $dx\wedge dy$, of the boundary
monomial $x^{\#\sigma_{1}}y^{\#\sigma_{2}}$ of the critical disk. With
$g_{m}=z^{m}(m_{2}x\partial_{x}-m_{1}y\partial_{y})$ the tropical-vertex generator
(Hamiltonian for $d\log x\wedge d\log y$), one has
$X_{k_{1},k_{2}}=(xy)^{-1}g_{(k_{1}+1,k_{2}+1)}$: the algebra is graded by an
integral-affine lattice, the origin shifted by the primitive form, and the
wall-crossings are the rational-power tropical vertex maps of
\cite[Lem.~4.20]{GKT}. The shift is the removal of the $\sigma_{12}$-insertion.
\end{proposition}

\begin{proof}
The first statement is \cite[Prop.~2.37(b)]{GKT} read against
Definition~\ref{def:diagonal}. The dictionary $X_{H}^{\mathrm{tor}}
=xy\,X_{H}^{\mathrm{aff}}$ follows from $\omega_{\mathrm{aff}}
=xy\,\omega_{\mathrm{tor}}$, and the explicit wall maps are
\cite[Lem.~4.20]{GKT}.
\end{proof}

\begin{proof}[Proof of Theorem~\ref{thm:trichotomy}]
The constraint \eqref{eq:dimconstraint} reads $sk_{1}+rk_{2}=w(J,\mathbf d)$ with
$k_{1},k_{2}\ge0$: for $w<0$ there are no solutions, for $w=0$ only $(0,0)$, and
for $w>0$ the $N$ solutions of Definition~\ref{def:diagonal}. A twist $(a,b)$
contributes $sa+rb-rs\le s(r{-}2)+r(s{-}2)-rs=rs(\chat_{W}-1)$ to the weight, with
equality at the maximal twist; descendents contribute $rs\,d_{j}>0$. Hence
positive-weight $d=0$ monomials exist if and only if $\chat_{W}>1$, and for
$\chat_{W}=1$ every $d=0$ monomial has $w\le0$ with equality exactly on products of
marginal twists; whenever such an equality monomial is wall-bearing, its unique
direction is the torus one. For the
existence of walls (not merely of positive-weight monomials) when
$\chat_{W}>1$: take $m$ interior markings of maximal twist $(r{-}2,s{-}2)$;
then $N=m-\lceil2m/r\rceil-\lceil2m/s\rceil+1\ge1$ once
$m(1-\tfrac2r-\tfrac2s)\ge2$, while the weight is $m\,rs(\chat_{W}-1)>0$, so
the resulting critical graphs carry wall summands; this is the maximal-twist
construction of \cite{note}. At $\chat_{W}=1$, those analogous multisets for
which $N\ge1$ are critical and marginal, realizing the torus wall-crossing; for $\chat_{W}<1$,
$\ell_{1}+\ell_{2}\le m\chat_{W}<m$ forces $N\le0$: no primary critical graphs
at all. The qualifier in the theorem is necessary: acting on a lower divisor
coefficient multiplies deformation monomials, so nonlinear transport can
contribute to a nonpositive-weight target even though no generator is supported there.
\end{proof}

\begin{remark}[The marginal torus direction]
\label{rem:torus}
At a wall-bearing marginal monomial ($w=0$ and $N\ge1$; necessarily
$r(J)=s(J)=0$ and $N=1$) the critical graph exists (for $x^{4}+y^{4}$ it
is precisely the disk of \cite[Fig.~9, Ex.~4.9]{GKTsurvey}), but its Lie
bidegree is $(0,0)$, the direction $x\partial_{x}-y\partial_{y}$ excluded from
$\gfrak^{r,s}$ in \cite[(4.15),(4.18)]{GKT} for nilpotency of the
$(x,y)$-filtration. The exclusion is an artifact: with coefficients in the ideal
generated by the deformation variables, the exponential converges adically, and the
proofs of \cite[Thm.~5.8]{GKT} (whose jump sets $Q_{0}(A,\mathbf d)$ impose only
\eqref{eq:dimconstraint}, admitting $(0,0)$) and of \cite[Cor.~5.13]{GKT} (whose
realizability step moves the two adjacent balanced coefficients by
$h\cdot((-1)^{k_{2}}s\,k_{1}(\Lambda),(-1)^{k_{2}+1}r\,k_{2}(\Lambda))$, at a
marginal diagonal exactly the antisymmetric torus motion) use the torus summands.
We therefore work with the enlarged algebra
\[
\widetilde\gfrak^{\,r,s}=\gfrak^{r,s}\ \oplus\
\bigoplus_{\substack{\varnothing\ne J,\ w(J,\mathbf d)=0,\ N(J,\mathbf d)\ge1\\
r(J)=s(J)=0}}
\QQ\,u_{J,\mathbf d}\,(x\partial_{x}-y\partial_{y}),
\]
whose new summands are abelian among themselves and act diagonally on the graded
pieces. This is a Lie algebra: for $D=X_{0,0}$,
$[uD,vX_{a,b}]=(a-b)uvX_{a,b}$ and $[uD,vD]=0$; since a marginal coefficient
$u$ has weight and residues zero, $uv$ satisfies the same critical-generator
constraints as $v$. In the primary algebra nothing is added when $\chat_{W}<1$; for the parabolic pairs the primary
group is purely torus, matching the elliptic wall-crossing group of Maher
\cite[\S5.2.2]{Maher}; for $\chat_{W}>1$ the torus summands sit over the marginal
submonomials. All statements below are over $\widetilde\gfrak^{\,r,s}$ and its
group $\widetilde G^{r,s}=\exp(\widetilde\gfrak^{\,r,s})$; the next lemma
supplies the torsor properties over the enlargement.
\end{remark}

\begin{lemma}[The enlarged torsor]
\label{lem:enlarged}
The action of $\widetilde G^{r,s}$ on potentials preserves $\Om(I)$ and is
\emph{free} and transitive on it, so that $\Om(I)$ is a torsor in the strict sense,
and the jump automorphisms of \cite[Thm.~5.8]{GKT} lie in $\widetilde G^{r,s}$.
\end{lemma}

\begin{proof}
Membership of the jumps: the jump sets $Q_{0}(A,\mathbf d)$
in the proof of \cite[Thm.~5.8]{GKT} impose only \eqref{eq:dimconstraint}, so
the jump elements are sums of pro-nilpotent and marginal torus summands.

Preservation and transitivity can be proved without inferring a full
exponential action from a single leading coefficient change. The
integration-by-parts argument of \cite[Thm.~4.26, Step~1]{GKT}, which proves
that the group preserves $\Upsilon_{r,s}(I)$, uses the dimension constraint and
the divergence-free Hamiltonian form of the generators. It applies unchanged
to $X_{0,0}=x\partial_x-y\partial_y$; its coefficient has positive deformation
degree, so the exponential converges in the deformation filtration. It also
preserves the chamber-index normalization: for a primary singleton, residue
zero forces twist $(0,0)$ and hence weight $-rs$, so a primary marginal torus
monomial has $|J|\ge2$ and acts trivially on the zero- and one-marking
normalized terms; a descendent singleton torus monomial does not alter their
$\mathbf d=0$ normalization (the normalizations proper,
\cite[Def.~3.46(1)]{GKT}, constrain only $\nu_{\Gamma,\mathbf 0}$; the
remaining conditions there are of $\Ascr$-form and are preserved by the
integration-by-parts argument). The
triangular construction of \cite[Thm.~4.26, Step~3]{GKT} then proves
transitivity. For the new degree-zero vector field, refine its weighted-degree
induction by deformation divisibility: a correction
$u_{J,\mathbf d}X_{0,0}$ changes its target diagonal at leading order and only
proper multiples thereafter. This secondary induction terminates modulo every
finite divisor-closed coefficient cutoff and passes to the coefficientwise
completion. The only new case is a marginal
diagonal: it has two balanced coefficients and its allowed difference is
exactly the one-dimensional antisymmetric direction
$r e_1-s e_0=X_{0,0}(W_0)$: a marginal diagonal is NS-visible
($r(J)=s(J)=0$), so its single $\Ascr$-condition is one nontrivial linear
condition on the two coefficients, and the torus vector spans its kernel (the
period computation in the proof of Theorem~\ref{thm:effective}); hence the
same induction closes. Finally
$\Om(I)=\Upsilon_{r,s}(I)$ by \cite[Cor.~5.13]{GKT}; hence the enlarged action
preserves and is transitive on the realizable indices.

Freeness: let $g=\exp(v)$ with $v\ne0$, and let $u$ be minimal in the
divisibility order among the monomials in the support of $v$. The change of
the $u$-diagonal of any potential is exactly $v_{u}(W_{0})$, $W_{0}=x^{r}+y^{s}$: contributions of
other summands, or of higher exponential terms, carry monomials properly
containing a support monomial, which cannot equal $u$ by minimality. And
$v_{u}(W_{0})\ne0$. Indeed, on the diagonal belonging to $u$, the $N$
pro-nilpotent generators act by the adjacent vectors
$a_{p}e_{p}-b_{p}e_{p-1}$ with $a_{p},b_{p}>0$. These vectors are linearly
independent: reading a relation successively in the
$e_{N},e_{N-1},\ldots,e_{1}$ coordinates kills all its coefficients. On a
marginal diagonal there is instead the single nonzero torus vector
$r e_{1}-s e_{0}$. Thus no nonzero linear combination of the summands in
$v_u$ can cancel. (Individually, a pro-nilpotent summand contributes
$r(k_{2}{+}1)x^{k_{1}+r}y^{k_{2}}-s(k_{1}{+}1)x^{k_{1}}y^{k_{2}+s}$, while a
torus summand contributes $r\,x^{r}-s\,y^{s}$.) So $g$ moves every potential, which is freeness and not
merely faithfulness. The distinction matters downstream: uniqueness of the
transition element, used in Theorem~\ref{thm:presentation}(i), needs trivial
stabilizers, and faithfulness alone does not supply them for a nonabelian group.
The argument of \cite[Thm.~4.26, Step~2]{GKT} likewise establishes freeness,
though it is stated there as faithfulness. All uses below are through this lemma.
\end{proof}

\section{The scattering presentation}
\label{sec:presentation}

Throughout, a \emph{normal-form family} is a family of simple canonical
homotopies supplied by \cite[Lem.~3.39]{GKT}: for each subset $A$ of the
internal labels there is one time $t_{A}\in(0,1)$, all times distinct, and
every jump of a critical graph with internal marking set $A$ occurs at
$t_{A}$. The times are auxiliary choices (Remark~\ref{rem:orderinggap}).

\begin{proof}[Proof of Theorem~\ref{thm:presentation}(i)]
Let $(\mathfrak s_{t})_{t\in[0,1]}$ be a family of simple canonical homotopies
supplied by \cite[Lem.~3.39]{GKT}, with $\mathfrak s_{1}=\mathfrak s_{0}$ and jump times
$t_{1}<\dots<t_{M}$ and jump elements $g_{i}\in\widetilde G^{r,s}$, each a sum over
the critical graphs with internal markings $A_{i}$, with coefficients, up to the
explicit signs $(-1)^{|A_{i}|+k_{2}}$ ($k_{2}$ the Lie bidegree) of the display
of $g_{i}$ there, the zero
counts $\Delta(\Lambda)$ of the homotopy on the critical stratum
(\cite[proof of Thm.~5.8, eqns.~(5.33)--(5.37)]{GKT}). By that proof the potential
$W^{\mathfrak s_{t}}$ is constant between jumps and transforms by $g_{i}$ across $t_{i}$;
around the loop, $W^{\mathfrak s_{0}}=(g_{M}\circ\cdots\circ g_{1})(W^{\mathfrak s_{0}})$, and
freeness (Lemma~\ref{lem:enlarged}) gives $g_{M}\circ\cdots\circ
g_{1}=\mathrm{id}$. For two such paths with fixed endpoints, their total products
$g,g'$ satisfy $g'g^{-1}(W^{\mathfrak s_{1}})=W^{\mathfrak s_{1}}$, hence $g'=g$.
\end{proof}

\begin{proof}[Proof of Theorem~\ref{thm:presentation}(ii)]
Let $\mathfrak m$ be the deformation ideal and first work in the primary
quotient
\[
 F^{p}\widetilde\gfrak^{\,r,s}_{0}
 :=\operatorname{span}\{u_{J,\mathbf0}X:\lvert J\rvert\ge p\},
 \qquad
 L_{<K}:=\widetilde\gfrak^{\,r,s}_{0}/F^{K}.
\]
It is a finite-dimensional nilpotent Lie algebra and
$[F^{p},F^{q}]\subseteq F^{p+q}$ (here and below $F^\bullet$ also denotes the
induced filtration on the quotient). For a rational boundary ray $\rho$ let
$L_{\rho}$ be the span of the generators $uX_{k_{1},k_{2}}$ for which
$(k_{1}{+}1,k_{2}{+}1)\in\rho$, placing a marginal generator
$u(x\partial_x-y\partial_y)$ in $\rho=\RR_{\ge0}(1,1)$. Then
$L_{<K}=\bigoplus_{\rho}L_{\rho}$, with only finitely many summands.

Each $L_{\rho}$ is abelian. Indeed,
\[
[X_{k_{1},k_{2}},X_{k'_{1},k'_{2}}]
=\det\bigl((k'_{1}{+}1,k'_{2}{+}1),(k_{1}{+}1,k_{2}{+}1)\bigr)
X_{k_{1}+k'_{1},\,k_{2}+k'_{2}},
\]
so the bracket vanishes when the shifted directions are parallel; on the
$(1,1)$-ray the marginal vector field also commutes with every $X_{q,q}$.
The deformation coefficient of a marginal generator has positive
$\mathfrak m$-degree, so it belongs to the same nilpotent filtration.

Fix any slope order. The multiplication map
\[
\prod_{\rho}^{\longrightarrow}\exp(L_{\rho})
\longrightarrow \exp(L_{<K})
\]
is bijective. Modulo $F^{2}$ this is the direct-sum decomposition above.
Suppose the factors have been found through $F^{n}$. The logarithm of the
residual in $F^{n}/F^{n+1}$ decomposes uniquely by $\rho$; insert these
components into the corresponding factors. Commuting a degree-$n$ correction
past a previously chosen positive-degree term changes the product only in
$F^{n+1}$. This proves existence by induction. At the first filtration degree
where two ordered factorizations differ, the same direct-sum decomposition
forces all differences to vanish, proving uniqueness. The construction is
compatible as $K$ varies.

For descendents, perform the same construction after quotienting to any finite
divisor-closed set of deformation monomials. Its complement is upward closed,
so this is a Lie-algebra quotient, and the compatible finite constructions give
the asserted coefficientwise presentation.
\end{proof}

\begin{remark}
\cite[Thm.~5.8]{GKT} produces the total jump in a homotopy-time-ordered
factorization, one factor per subset of internal markings, mixing wall directions;
Lemma 3.39 of \cite{GKT} shows the jump times and their order are auxiliary
choices. The boundary-ray presentation is a canonical reorganization by the
\emph{shifted} directions forced by Proposition~\ref{prop:shift}, not by the
additive Lie rays $(k_{1},k_{2})$. The affine shift does not preserve rays; this
filtered distinction is what makes the canonical diagram of
Theorem~\ref{thm:canonical} geometrically meaningful.
\end{remark}

\section{Per-diagonal structure: the recursion transverse to the walls}
\label{sec:diagonal}

Two structural facts from \cite{GKT} organize this section. The open
topological recursion factorizes through a
cohomological-field-theory node splitting: in
\cite[Lem.~6.4]{GKT} the distinguished insertion splits off a closed extended
invariant and continues as an open piece carrying the socle-dual twist
$(a,b)\mapsto(r{-}2{-}a,s{-}2{-}b)$ of the Jacobi ring
$\CC[x,y]/(x^{r-1},y^{s-1})$: the node is the residue pairing. And the
$\Gamma$-weights in the invariant combinations $\Ascr$ are the periods of
$x^{r}+y^{s}$: $\int_{0}^{\infty}e^{-x^{r}}x^{k}dx=\tfrac1r\Gamma(\tfrac{1+k}r)$,
so the weight attached to a boundary profile is the period normalization of the
primitive form. The good-basis period moments themselves are explicit:

\begin{lemma}[Periods]
\label{lem:periods}
For the good-basis cycles $\Xi_{a}$ of $x^{r}$
($\int_{\Xi_{a}}x^{a'}e^{x^{r}/\hbar}=\delta_{aa'}$, $0\le a,a'\le r-2$), the
moments $P_{a}(p)=\int_{\Xi_{a}}x^{p}e^{x^{r}/\hbar}$ satisfy
$P_{a}(p)=\delta_{ap}$ for $p\le r{-}2$, $P_{a}(r{-}1)=0$, and
$P_{a}(p)=-\tfrac{\hbar(p-r+1)}{r}P_{a}(p-r)$.
\end{lemma}

\begin{proof}
Total derivatives integrate to zero on $\Xi_{a}$:
$0=\int_{\Xi_{a}}\tfrac{d}{dx}(x^{q}e^{x^{r}/\hbar})
=q\,P_{a}(q{-}1)+\tfrac r\hbar P_{a}(q{+}r{-}1)$.
\end{proof}

\begin{proof}[Proof of Theorem~\ref{thm:effective}]
\emph{The case $N=0$.}
If $N=0$, there is no wall direction and the nonzero $h=1$ weight in
$\Ascr$ determines the sole coefficient modulo lower monomials, so the claim is
immediate.

\emph{Wall-bearing diagonals.}
Fix now a wall-bearing diagonal $(J,\mathbf d)$, of positive or
marginal weight, with coefficients $\nu_{p}$, $0\le p\le N$, as in
Definition~\ref{def:diagonal}. The wall element of
$\Lambda_{J,p}$ acts on $W_{0}=x^{r}+y^{s}$ by
\begin{equation}
\label{eq:wallvector}
u_{J,\mathbf d}X_{k_{1},k_{2}}(W_{0})
=u_{J,\mathbf d}\bigl(r(k_{2}{+}1)\,x^{k_{1}+r}y^{k_{2}}
-s(k_{1}{+}1)\,x^{k_{1}}y^{k_{2}+s}\bigr),
\end{equation}
which lands on the two adjacent balanced points of the same diagonal: in the basis
$e_{0},\dots,e_{N}$ of coefficients, the wall $p$ moves along the jump vector
$w_{p}=a_{p}e_{p}-b_{p}e_{p-1}$ (not to be confused with the scalar Euler
weight $w(J,\mathbf d)$) with $a_{p}=r\,k_{2}(\Lambda_{J,p})>0$,
$b_{p}=s\,k_{1}(\Lambda_{J,p})>0$; here and below $k_{i}(\Lambda)$ denotes the
\emph{boundary} degrees of $\Lambda$, which exceed the Lie bidegree
$(k_{1},k_{2})$ of the display by $(1,1)$, so that $r(k_{2}{+}1)=r\,k_{2}(\Lambda)$
and $s(k_{1}{+}1)=s\,k_{1}(\Lambda)$, both positive. (These are exactly the geometric jump
coefficients of \cite[(5.51)--(5.53)]{GKT}.) The $w_{p}$ are triangular, hence
independent, and each is annihilated by the period functional: by
Lemma~\ref{lem:periods}, $P_{a}(k_{1}{+}r)P_{b}(k_{2})
=-\tfrac{\hbar(k_{1}+1)}{r}P_{a}(k_{1})P_{b}(k_{2})$ cancels against the second
term of \eqref{eq:wallvector}. \emph{NS-visible diagonals.}
On an NS-visible diagonal ($r(J)\le r-2$ and
$s(J)\le s-2$) the period functional is nonzero on each $e_{p}$, so its kernel has
dimension $N$ and coincides with $\mathrm{span}(w_{p})$. By
\cite[Thm.~4.9]{GKT} the coefficient of $u_{J,\mathbf d}$ in
$\int_{\Xi_{(r(J),s(J))}}e^{W^{\nu}/\hbar}\,dx\wedge dy$ equals
$(-1)^{|J|}(-\hbar)^{-d(J,\mathbf d)-2}\Ascr(J,\mathbf d,\nu)$, whose $h=1$ part
($h$ the number of parts in the ordered-partition expansion of $\Ascr$,
\cite[Thm.~4.9]{GKT}) is
the period-weighted sum $\sum_{p}\omega_{p}\nu_{p}$ over the diagonal and whose
$h\ge2$ parts involve only strictly smaller monomials. The chamber-index
conditions \cite[Def.~3.46]{GKT} prescribe the value of $\Ascr$.

\emph{Invisible diagonals.}
On invisible
diagonals ($r(J)=r-1$ or $s(J)=s-1$) the good-basis periods vanish identically
on the diagonal (Lemma~\ref{lem:periods}), but the $h=1$ weights of $\Ascr$
remain nonzero there (their arguments $(1+K)/r$ are positive integers,
where $\Gamma$ does not vanish), and $\mathrm{span}(w_{p})$ lies in the
kernel of the weighted functional modulo strictly smaller monomials, by the
boundary-condition independence of $\Ascr$ \cite[Thm.~5.4]{GKT}; so the same
one-linear-condition structure holds, imposed directly by the conditions,
which by \cite[Cor.~5.13]{GKT} characterize $\Om(I)$ completely. Ordering monomials by degree and descendents determines every
invariant, up to the span of the $w_{p}$ (the torsor directions), proving the
theorem. The recursion \cite[Thm.~0.5]{GKT} is the same system of conditions
organized by descendent reduction, whence the last statement.
\end{proof}

\begin{remark}[The first computed tower]
\label{rem:tower}
Call the family of diagonals over the powers of one deformation parameter its
\emph{tower}. For $x^{5}+y^{5}$ the conditions produce, order by order over the
$t_{3,3}$-monomials (all verified by exact arithmetic; \S\ref{sec:cert}): the singleton normalization
$\langle\tau_{0}^{(3,3)}\sigma_{1}^{3}\sigma_{2}^{3}\sigma_{12}\rangle=1$; the
first-wall relation \eqref{eq:firstwallrel} at $t_{3,3}^{2}$; the absolute value
$\tfrac15$ of \eqref{eq:newinvariants} at $t_{3,3}^{3}$ (a diagonal with $N=0$,
hence no walls, and $r(J)=4=r{-}1$, hence invisible to the good-basis periods,
so the condition that determines it is not a period condition); and at $t_{3,3}^{4}$
a relation quadratic in the first-wall invariants, giving
$\nu^{\min}_{(2,7)}=-\tfrac{12}{25}$ in the normal form of
Theorem~\ref{thm:canonical}. The descendent seed at $d=1$ gives
$\langle\tau_{1}^{(3,3)}\sigma_{1}^{8}\sigma_{2}^{3}\sigma_{12}\rangle
+\langle\tau_{1}^{(3,3)}\sigma_{1}^{3}\sigma_{2}^{8}\sigma_{12}\rangle
=-\tfrac54$.
\end{remark}

\begin{remark}[The seed values, and {\cite[Ex.~5.5]{GKTsurvey}}]
\label{rem:ex55}
Imposing \cite[Thm.~0.3]{GKT} on the oscillatory expansion of the potential
directly --- the mechanism by which \cite[Thm.~0.2]{GKT} is proved, so that no
combinatorial convention enters --- forces, at every two-insertion seed diagonal
$\{(a_{1},b_{1}),(a_{2},b_{2})\}$ with $a_{1}{+}a_{2}=r$, $b_{1}{+}b_{2}=s$,
$\mathbf d=0$, the relation \eqref{eq:seed}: the diagonal supports the boundary
profiles $(r,0)$ and $(0,s)$ only, $d(J,0)<0$ so the mirror theorem admits no term
there, and the linear and quadratic insertions contribute
$\nu_{(r,0)}/r+\nu_{(0,s)}/s$ and $1/(rs)$ respectively
(Lemma~\ref{lem:periods}; verified symbolically for $r=s=4,5$ and
$(r,s)=(3,5),(4,6),(3,7)$, \S\ref{sec:cert}).
At $r=s=4$ this gives $-\tfrac14$ for the sum considered in
\cite[Ex.~5.5]{GKTsurvey}, where $-\tfrac18$ is stated; the discrepancy traces to
the $h=2$ partition count in the evaluation of $\Ascr$ there, which
\cite[Thm.~4.9]{GKT} fixes as the count of ordered partitions of the labeled
marking set with $1/h!$. The same count is pinned independently by the proven
\cite[Thm.~5.4]{GKT}: at $r=s=5$, $\Ascr(\{(1,1),(1,1)\},0)$ equals the closed
three-point value $\langle\tau_{0}^{(1,1)}{}^{3}\rangle=1$, which the labeled
count reproduces and the halved count does not (\S\ref{sec:cert}). On the marginal seed diagonals the two
coefficients in \eqref{eq:seed} are moved (antisymmetrically) only by the torus
direction of Remark~\ref{rem:torus}; the symmetric gauge value is
$\nu_{(r,0)}=\nu_{(0,s)}=-\tfrac1{2r}$ at $r=s$. The discrepancy concerns
only the worked example; no theorem of \cite{GKT} or \cite{GKTsurvey} is
affected.
\end{remark}

\section{Normal forms and the canonical diagram $\Dscr^{r,s}$}
\label{sec:canonical}

\begin{proof}[Proof of Theorem~\ref{thm:canonical}]
Work in a primary quotient modulo $F^K$, or more generally on a finite
divisor-closed coefficient set as in Theorem~\ref{thm:presentation}(ii).
For existence and uniqueness of $\nu^{\min}$, order its monomials by total
degree. At each step, on each wall-bearing diagonal the span of the wall
vectors $w_{p}=a_{p}e_{p}-b_{p}e_{p-1}$ ($a_{p},b_{p}>0$) is complementary to
$\QQ e_{0}$: a combination $\sum c_{p}w_{p}\in\QQ e_{0}$ forces, reading the
$e_{N},e_{N-1},\dots$ components in turn, $c_{N}=\dots=c_{1}=0$. Hence there is a
unique group correction at this order making all $\nu_{p}$, $p\ge1$, vanish, and
it does not disturb lower orders; on marginal diagonals the torus direction plays
the role of the single wall. The resulting index is enumerative because the group
preserves $\Om(I)$ (Lemma~\ref{lem:enlarged}). Uniqueness follows from
freeness (Lemma~\ref{lem:enlarged}) plus the same triangularity: a nontrivial leading correction moves some
$\nu_{p\ge1}$. The construction for $\nu^{\max}$ is symmetric (complement of
$\QQ e_{N}$). The boundary-slope-ordered factorization of the unique
$g^{\min\to\max}$ (Theorem~\ref{thm:presentation}(ii)) is then canonical, and its
wall functions are computed order by order from the normal forms. Uniqueness
makes these constructions compatible under restriction, giving the primary
completion and the coefficientwise descendent extension.
\end{proof}

\begin{example}[The first walls of $\Dscr^{5,5}$]
\label{ex:walls}
At order two in the deformation variables the three wall-carrying diagonals of
$x^{5}+y^{5}$ have balanced profiles $(0,6),(5,1)$; $(6,0),(1,5)$; and
$(1,6),(6,1)$ respectively, with invariant relations and normal forms (\S\ref{sec:cert}):
{\small
\[
\renewcommand{\arraystretch}{1.35}
\begin{array}{lll}
\toprule
\text{monomial} & \text{relation} & \text{normal forms}\\
\midrule
t_{2,3}t_{3,3} &
\tfrac25\nu_{(0,6)}+\tfrac15\nu_{(5,1)}=-\tfrac2{25} &
\nu^{\min}_{(0,6)}=-\tfrac15,\ \nu^{\max}_{(5,1)}=-\tfrac25\\
t_{3,2}t_{3,3} & \text{the image under $x\leftrightarrow y$} & \text{by symmetry}\\
t_{3,3}^{2} &
\tfrac25\bigl(\nu_{(1,6)}+\nu_{(6,1)}\bigr)=-\tfrac4{25} &
\nu^{\min}_{(1,6)}=\nu^{\max}_{(6,1)}=-\tfrac25\\
\bottomrule
\end{array}
\]
}
Solving $W^{\nu^{\max}}-W^{\nu^{\min}}=c\,\cdot u\,X_{k_{1},k_{2}}(W_{0})$ on each
diagonal is an overdetermined system (two monomial components, one
coefficient) whose consistency is forced by the relations above and constitutes
a check of the whole mechanism; it yields the wall functions of
Theorem~\ref{thm:canonical}. The wall $t_{3,3}^{2}X_{1,1}$ is the first symmetric
pro-nilpotent primary wall: only the sum \eqref{eq:firstwallrel} is invariant,
the difference jumps.
\end{example}

\section{The closed vertex data for $x^{5}+y^{5}$, and the verified recursion}
\label{sec:x5y5}

The determination of Theorem~\ref{thm:effective} consumes closed extended FJRW
invariants. For the first symmetric hyperbolic Fermat pair $x^{5}+y^{5}$
($\chat_{W}=\tfrac65>1$)
we compute this input exactly.

\begin{example}[Closed vertex weights for $x^{5}+y^{5}$]
\label{ex:x5y5}
By Sebastiani--Thom the state space and cohomological field theory of the Fermat
sum factor as $A_{4}(x)\otimes A_{4}(y)$, the tensor square of the $5$-spin
($A_{4}$) theory; a closed invariant
$\langle\prod_{i}\tau_{d_{i}}^{(a_{i},b_{i})}\rangle$ is the tensor correlator
$\int_{\Mbar_{0,n}}W^{x}(\mathbf a)\,W^{y}(\mathbf b)\,\psi^{\mathbf d}$ with
$W^{x},W^{y}$ the two Witten classes, of degrees $\varepsilon_{1}=(\sum a_{i}-3)/5$ and
$\varepsilon_{2}=(\sum b_{i}-3)/5$. We compute these exactly (primitive-form normalization,
sign $(-1)^{\varepsilon_{1}+\varepsilon_{2}}$) by three independent, mutually consistent routes:
\begin{itemize}\itemsep2pt
\item the single-variable $A_{4}$ invariants from the Saito primitive form (flat
coordinates, three-point residues, string/dilaton/topological-recursion
descendents), validated against the $A_{2}$ normalization
$\langle(\tau_{0}^{1})^{4}\rangle=\tfrac13$, WDVV for $A_{2},A_{3},A_{4}$, and the
$r=2$ Witten--Kontsevich numbers;
\item the factorizing tensor correlators ($\varepsilon_{1}=0$ or $\varepsilon_{2}=0$), where one
Witten class is the topological unit;
\item the fully mixed correlators with $\min(\varepsilon_{1},\varepsilon_{2})=1$, by reconstructing
the divisor $\mathcal W=\sum_{U}c_{U}\delta_{U}$ on $\Mbar_{0,n}$ from its node-splitting
restrictions and pairing with the other class, using a self-contained genus-zero
tautological intersection calculus (splitting recursion,
$\delta_{S}^{2}=-\psi'-\psi''$), verified against Witten--Kontsevich and the
Petersen intersection form of $\Mbar_{0,5}=\mathrm{dP}_{5}$; each reconstructed
class passes the independent check $\langle \mathcal W\,\psi^{n-4}\rangle
=\langle\tau_{1}^{a_{i}}\cdots\rangle_{A_{4}}$.
\end{itemize}
Representative values (denominators $5^{\,n-3}$; the socle pairing is
$\langle\tau_{0}^{(0,0)}\tau_{0}^{(0,0)}\tau_{0}^{(3,3)}\rangle=1$):
\[
\renewcommand{\arraystretch}{1.4}
\begin{array}{lll}
\langle(\tau_{0}^{(2,1)})^{3}\tau_{0}^{(2,0)}\rangle=-\tfrac25, &
\langle\tau_{0}^{(1,3)}\tau_{0}^{(1,0)}(\tau_{0}^{(3,0)})^{2}\rangle=-\tfrac15
& (\varepsilon_{1},\varepsilon_{2})=(1,0),\\
\langle(\tau_{0}^{(2,1)})^{2}\tau_{0}^{(2,2)}(\tau_{0}^{(1,2)})^{2}\rangle
=\tfrac{4}{25}, &
\langle(\tau_{0}^{(2,2)})^{3}(\tau_{0}^{(1,1)})^{2}\rangle=\tfrac{2}{25}
& (1,1),\\
\langle(\tau_{0}^{(2,3)})^{2}\tau_{0}^{(1,3)}\tau_{0}^{(1,2)}
(\tau_{0}^{(1,1)})^{2}\rangle=-\tfrac{4}{125}, &
\langle(\tau_{0}^{(2,3)})^{2}(\tau_{0}^{(1,3)})^{3}\tau_{0}^{(1,2)}
\tau_{0}^{(0,1)}\rangle=\tfrac{6}{625}
& (1,2),\ (1,3).
\end{array}
\]
The smallest case with both classes of codimension two, $n=7$, is computed by
the codimension-two stratum intersection form: e.g.\
$\langle\tau_{0}^{(3,1)}\tau_{0}^{(2,1)}(\tau_{0}^{(2,2)})^{3}\tau_{0}^{(1,2)}
\tau_{0}^{(1,3)}\rangle=\tfrac{34}{625}$, invariant under relabelling. All
values are verified by the exact computations of \S\ref{sec:cert}.
\end{example}

As independent evidence for the entire framework, the open topological recursion
of \cite[Thm.~0.5]{GKT} is confirmed numerically for $x^{5}+y^{5}$: writing the
canonical combinations $\Ascr(A,\mathbf d)$ as closed extended invariants
(\cite[Thm.~5.4]{GKT}) and evaluating both sides,
an exhaustive check covers every instance with insertions in the narrow
range, and a second, independent computation removes the restriction by
computing every invariant, Neveu--Schwarz and extended (Ramond,
$\mathrm{tw}=-1$) alike, through Witten's $r$-spin (Chiodo) class in SageMath's
\texttt{admcycles}: the recursion holds in all $72$ instances with $|J|=2$
and all $105$ with $|J|=3$ (\S\ref{sec:cert}). As a
by-product, the \texttt{admcycles} route independently reproduces every closed
invariant of
Example~\ref{ex:x5y5}, including $\tfrac{34}{625}$ and the descendent values.

\section{The geometric picture}
\label{sec:geometric}

The results above define the canonical diagram $\Dscr^{r,s}$ and show that its
consistency structure determines the invariants. What they do not provide is a
\emph{position space}: an intrinsic domain in which the walls are rays and the
invariants are computed by broken lines, as for the Fano Landau--Ginzburg models
of \cite{Gross10,CPS}. This section builds as much of that picture as is a theorem
and isolates what is not.

Three words of warning on vocabulary. We reserve
\emph{base} for the base $B\Gamma$ of a graded graph \cite[Def.~3.5]{GKT}, an
unrelated notion carrying the inductive structure of the boundary conditions, and
call the domain below the \emph{direction sector}; and \emph{consistency} here
always means the path-independence of Theorem~\ref{thm:presentation}(i) in
boundary-condition space.  The product $g^{\min\to\max}$ is the transport of a
boundary-to-boundary sweep of the sector, whose two ends carry $\nu^{\min}$ and
$\nu^{\max}$; it is not loop holonomy
(Theorem~\ref{thm:positionobstruction7}).  Likewise \emph{chamber index} is
\cite{GKT}'s name for a system of invariants and does not refer to a chamber
of the direction sector.

The order below puts what is proved first. \S\ref{ssec:sector} places the walls,
\S\ref{ssec:transport} proves the transport representation, its ordered-chain
inversion and its corrected Euler-cascade class, and the sense in which the
theory exceeds its own tropicalization; \S\ref{ssec:realization} realizes every
finite coefficient window by staged homotopies, makes the realizations
coefficientwise compatible, and proves the exact obstruction that forces the
staging; and \S\ref{ssec:missing} proves the projective coherent-sector theorem
and the resulting position-space obstruction.

\subsection{The sector and its rays}
\label{ssec:sector}
\begin{definition}
\label{def:sector}
The \emph{boundary-charge lattice} is
$N_{\partial}=\ZZ e_{1}\oplus\ZZ e_{2}$, where $e_i$ records one boundary
insertion of type $\sigma_i$.  The \emph{direction sector} is its effective cone
$C_{\partial}=\RR_{\ge0}e_1+\RR_{\ge0}e_2$, with the $(1,1)$-shifted
integral-affine structure of Proposition~\ref{prop:shift}; its two sides are the
rays of the two boundary types. A wall of $\Dscr^{r,s}$ is placed on the
ray $\RR_{\ge0}(k_{1}{+}1,k_{2}{+}1)$ of its critical graph's boundary monomial,
and a chamber of a finite primary truncation is a connected component of the
complement of its rays. For the completed descendent diagram, a \emph{slope cut}
is a direction together with a choice of side when the direction itself carries
a wall. The \emph{base cut} is the side of the sector along the
$\sigma_{1}$-axis ray, and the transport $\Psi(q)$ to a cut $q$ is the
slope-ordered product of the wall
factors between the base cut and $q$, defined on each finite divisor-closed
coefficient set and
then coefficientwise; only the composites $\Psi(q)\Psi(q')^{-1}$ enter below,
so the choice of base is a convention. For $u,v\in\ZZ^{2}$ write
$\operatorname{cross}(u,v)=u_{1}v_{2}-u_{2}v_{1}$.
\end{definition}

Thus the quadrant, its lattice, and its axes come from the open FJRW boundary
data rather than from an auxiliary drawing. Only its projectivization
$\mathbb P(C_{\partial})\cong[0,\infty]$, the ordered slope space, is used by
transport. The placement is not a convention but forced, and the same computation shows why
the theory looks tropical at all: the $(1,1)$-shift of
Proposition~\ref{prop:shift}, which is nothing but the removal of the
$\sigma_{12}$-insertion, is exactly what makes bending along a transported
monomial monotone in slope, hence what makes broken lines behave like broken
lines here (Theorem~\ref{thm:nobackscatter}).


\begin{proposition}[Forced placement]
\label{prop:placement}
The wall element $X_{k_{1},k_{2}}$ fixes exactly the monomials proportional to
its boundary-monomial direction $(k_{1}{+}1,k_{2}{+}1)$. Hence in any geometric
realization the wall of a critical graph must lie on the ray of its boundary
monomial (the $(1,1)$-shifted ray), and a seed monomial is supported on the ray
of its own exponent. The direction sector of Definition~\ref{def:sector} is
thus forced, its sides the rays of the two boundary types; a full loop would have holonomy
$g^{\min\to\max}\ne\mathrm{id}$.
\end{proposition}

\begin{proof}
$X_{k_{1},k_{2}}(x^{a}y^{b})=\bigl((k_{2}{+}1)a-(k_{1}{+}1)b\bigr)
x^{a+k_{1}}y^{b+k_{2}}$, and the coefficient is the determinant of $(a,b)$
against $(k_{1}{+}1,k_{2}{+}1)$, vanishing precisely on the ray. The holonomy
statement is Theorem~\ref{thm:canonical}.
\end{proof}

\begin{proposition}[Interlacing and extreme support]
\label{prop:interlace}
On every diagonal the balanced monomials and the walls interlace in slope
(writing $\mathrm{bal}_{p}$ for the balanced monomial of $\Gamma_{J,p}$),
$\mathrm{bal}_{N}<\Lambda_{N}<\dots<\Lambda_{1}<\mathrm{bal}_{0}$, and the
normal forms of Theorem~\ref{thm:canonical} are the extreme-chamber support
conditions in every finite primary truncation: near one axis only the monomial at the far end of each diagonal
survives.
\end{proposition}

\begin{proof}
Immediate from the formulas of Definition~\ref{def:diagonal}: the wall
$\Lambda_{p}$ has boundary degree between those of $\mathrm{bal}_{p}$ and
$\mathrm{bal}_{p-1}$ in slope; the support statement is the normal-form
construction read through the interlacing.
\end{proof}



\subsection{The transport representation}
\label{ssec:transport}

\begin{definition}
\label{def:anchor}
The \emph{anchor cut} $q_m$ of a potential monomial $m$ is the one-sided cut
through its own direction, on the side toward the $(1,1)$-direction. If $m$ is
off every wall in a finite truncation, this is represented by the chamber
containing its direction; for $m$ on the $(1,1)$-direction either side gives the
same coefficient, and the exponent $(0,0)$ is anchored at the base cut. Its
\emph{anchored value} is its coefficient at $q_m$. For a
candidate assignment of seeds,
$T(q):=\sum_{m}\Psi(q)\Psi(q_m)^{-1}(\mathrm{seed}_{m})$ denotes the resulting
transport potential.
\end{definition}

Every bend along a transported monomial adds $\beta-(1,1)$ to the exponent,
$\beta$ the boundary direction on the crossed ray. Throughout, seed and
target are taken off the crossed rays: a seed on a bend ray has vanishing
bend coefficient, and an on-ray target uses the one-sided anchor cut toward
the $(1,1)$-direction, the convention of Theorem~\ref{thm:transport}. Call a
bending contribution to a coefficient \emph{foreign} if it originates from a
seed other than the coefficient's own.

\begin{theorem}[No backwards bending]
\label{thm:nobackscatter}
Along a path with $b$ bend factors from a seed with exponent $m_{\sigma}$
into a monomial $m$,
\[
b\,\operatorname{cross}(m,(1,1))
=\operatorname{cross}(m,m_{\sigma})+\sum_{f}\operatorname{cross}(m,\beta_{f}).
\]
Consequently no foreign bending reaches a monomial's anchor cut from the far
side of the monomial's own direction, the side away from the
$(1,1)$-direction: a monomial below the $(1,1)$-direction receives foreign
bending only along paths descending in slope, one above it only along paths
ascending, and monomials on the $(1,1)$-direction
receive no foreign contribution at all.
\end{theorem}

\begin{proof}
Summing the per-bend exponent motion gives
$m+b(1,1)=m_{\sigma}+\sum_{f}\beta_{f}$; the displayed identity is its cross
product with $m$. Transport is monotone, so every crossed ray lies strictly
between the seed's direction and $m$'s. For a path from below, seed direction
and all crossed rays are strictly shallower than $m$, so the right side is
negative, forcing $\operatorname{cross}(m,(1,1))<0$, i.e.\ $m$ strictly above
the $(1,1)$-direction; from above it is positive, forcing $m$ strictly below. A monomial on the
$(1,1)$-direction, where the cross product vanishes, excludes both.
\end{proof}

The primitive-form shift is thus precisely the monotone-bending mechanism.

Throughout, \emph{seed} means the initial datum of a broken line in the sense of
\cite{Gross10,CPS} and carries no cluster-theoretic meaning.

\begin{lemma}[Reduction]
\label{lem:reduction}
For the anchored-value assignment, $T=W$ at every cut if and only if, for
every coefficient, the foreign bending contributions cancel at its anchor cut.
\end{lemma}

\begin{proof}
Wall automorphisms are additive, so $W-T$ transports as a family; its
coefficients on any divisibility-minimal nonzero monomial are
cut-independent (all bending contributions come from proper divisors, where $W-T$
vanishes), and evaluating each coefficient at its anchor cut gives zero
exactly under the stated cancellation.
\end{proof}

\begin{theorem}[Linear order; singleton towers]
\label{thm:linearorder}
At linear order in a diagonal's own wall functions applied to the axis seeds
$x^{r},y^{s}$, the cancellation of
Lemma~\ref{lem:reduction} holds on every diagonal of every
$(r,s,J,\mathbf d)$: transporting into the cut at $\mathrm{bal}_{p}$ from below
crosses exactly the walls $\Lambda_{q}$, $q\ge p{+}1$, whose $x^{r}$-bends land
on $\mathrm{bal}_{q}$, never on the target coefficient, while from above the
crossed walls' $y^{s}$-bends land on $\mathrm{bal}_{q-1}$, $q\le p$, again off
the target: the interlacing pattern is
the no-backwards-bending property of broken lines in \cite{CPS}. Bends of
divisor seeds at foreign walls are outside this statement; Example~\ref{ex:farside}
shows they can arrive. For singleton
diagonals (one internal marking, any descendent) no cross-diagonal contribution exists,
the linearization is exact, and the transport representation holds with the
anchored-value assignment on every such tower, with wall functions the chain
$c_{q+1}=c_{q}a_{q}/b_{q+1}$ closed by the $\Ascr$-condition.
\end{theorem}

\begin{proof}
The crossing sets and landing indices are those of
Proposition~\ref{prop:interlace} together with the landing identities: by
\eqref{eq:wallvector}, $X_{\Lambda_{q}}(x^{r})$ is a nonzero rational
multiple of the monomial $\mathrm{bal}_{q}$ and $X_{\Lambda_{q}}(y^{s})$ of
$\mathrm{bal}_{q-1}$. For a singleton
monomial no proper divisor exists, so the only contributions are the linear
ones.
\end{proof}

Two computed examples below (Examples~\ref{ex:middleseeds}
and~\ref{ex:farside}) calibrate what the seeds are.

\begin{theorem}[Transport representation; $=$ Theorem~\ref{thm:transport}]
\label{thm:transport7}
With the anchor assignment of Definition~\ref{def:anchor}, the difference between
the potential's coefficient vector on a monomial and the transported contributions
of the seeds of its proper divisors is independent of the cut. Defining the
\emph{seed system} by these constants, every cut potential is the transport
\[
W(q)=\sum_{m}\Psi(q)\,\Psi(q_m)^{-1}(\mathrm{seed}_{m})
\]
of the seeds, and the seed system is the unique one with this property.
\end{theorem}

\begin{proof}[Proof of Theorem~\ref{thm:transport}]
Induct on the divisibility order of deformation monomials; the base is
$W_{0}=x^{r}+y^{s}$, whose two coefficients are seeds with value $1$ in every
cut. Fix a monomial $u=u_{J,\mathbf d}$ and assume the seeds are defined,
with the transport identity holding, on every proper divisor of $u$. Write
$D$ for the diagonal of $u$, so that $u_{D}=u$, and $E$ for the diagonal of a
proper divisor $u_{E}$ of $u$. Let
$F(q)$ denote the $u$-component of the sum of the transported lower seeds.

\emph{Step 1: the jump formula.}
Across any wall $\theta=\exp(\eta)$ the coefficients on $D$ jump by
\[
\Delta_{\theta}\bigl(W|_{D}\bigr)
=\sum_{j\ge1}\frac1{j!}
\sum_{\substack{\eta_{1},\dots,\eta_{j}\in\operatorname{supp}(\eta)\\
u_{E}\,u_{\eta_{1}}\cdots\,u_{\eta_{j}}=u_{D}}}
\bigl[\eta_{1}\cdots \eta_{j}\bigl(W|_{E}\bigr)\bigr]_{D},
\]
a universal expression in the wall function and the coefficients of $W$ on
proper divisors $E$ of $D$ and on $W_{0}$: any other source overshoots
$u$, and $D$'s own walls act on $D$'s coefficients only through $W_{0}$
(their monomial is $u$ itself), so the jump is exactly linear there.

\emph{Step 2: cancellation, and chamber-independence of the residual.}
The jump of $F$ across the same wall is the same universal expression
evaluated on the transported-lower-seed sum, whose divisor coefficients equal
$W$'s by the induction hypothesis and whose $W_{0}$-coefficients are the same
constants $1$. Hence the two jumps cancel in the difference and
$V:=W(q)|_{D}-F(q)$ is independent of $q$.

\emph{Step 3: the seed at $D$, and uniqueness.}
A $u$-term transports with unchanged
$u$-component (every wall action on it overshoots $u$), so placing
$\mathrm{seed}(u):=V$ at the assigned anchor cuts gives
$T|_{D}=V+F=W|_{D}$ at every cut (all sums are coefficientwise finite: only
finitely many divisor seeds and bend factors contribute to a fixed monomial),
completing the induction. Uniqueness is
immediate: $V$ is determined by $W$ and the lower seeds. Anchor assignments
affect only the higher-order terms of the transports, hence the seed values
at higher orders, never the existence or uniqueness of the decomposition; the
$(1,1)$-side assignment is the canonical choice, being the one compatible with
Theorem~\ref{thm:nobackscatter}.
\end{proof}

\begin{lemma}[Ordered bend expansion]
\label{lem:bendchain}
Work in a finite divisor-closed coefficient quotient and write
$z^a=x^{a_1}y^{a_2}$.  Suppose a transport path crosses, in chronological
order,
\[
 \Theta_j=\exp(\!\epsilon_j\eta_j),\qquad
 \eta_j=\sum_{\lambda\in L_j}c_\lambda u_\lambda X_{k_\lambda},
 \qquad \epsilon_j\in\{1,-1\}.
\]
For nonnegative integers $n_j$ and ordered choices
$\lambda_{j,1},\ldots,\lambda_{j,n_j}\in L_j$, read the decorations first by
increasing $j$ and put $a^h=a^{h-1}+k_{\lambda_h}$, starting from $a^0=a$.
Then
\begin{equation}
\label{eq:bendchain}
 (\Theta_m\circ\cdots\circ\Theta_1)(u_Ez^a)
 =\sum_\gamma
 \frac{\prod_h\epsilon_{j(h)}c_{\lambda_h}
  \det(a^{h-1},k_{\lambda_h}+(1,1))}
 {\prod_j n_j!}\,
 u_E\prod_hu_{\lambda_h}\,z^{a^{|\gamma|}},
\end{equation}
the finite sum over the surviving decorated bend chains $\gamma$.
\end{lemma}

\begin{proof}
Expand each exponential in powers of $\eta_j$ and distribute each power over
its summands.  Successive applications use
\[
 X_k(z^a)=\bigl((k_2+1)a_1-(k_1+1)a_2\bigr)z^{a+k}
 =\det(a,k+(1,1))z^{a+k}.
\]
The exponential supplies $1/n_j!$, the path supplies the displayed factor
order, and nilpotence makes the sum finite.  Generators on one boundary ray
commute, so repeated bends may equivalently be grouped as multisets with the
same factorials.  Notice that the expansion is by chains: transport is linear
in its initial monomial.  Rooted Lie trees arise only after an additional,
noncanonical BCH expansion of the ray factors.  \qedhere
\end{proof}

Index a potential slot by $i=(D,a)$, where $D$ is a deformation monomial and
$a$ an exponent on its diagonal; let $q_i$ be its anchor cut.  If
$j=(E,b)$, define the strictly divisibility-triangular transport kernel
\begin{equation}
\label{eq:transportkernel}
 K_{ij}=
 [u_Dz^a]\,\Psi(q_i)\Psi(q_j)^{-1}(u_Ez^b)
 \quad\text{when }E\mid D, E\ne D,
 \qquad K_{ij}=0\quad\text{otherwise}.
\end{equation}
Lemma~\ref{lem:bendchain} is an explicit formula for $K_{ij}$: retain the
chains whose deformation labels multiply $u_E$ to $u_D$ and whose exponent
increments carry $b$ to $a$.

\begin{theorem}[Triangular seed inversion]
\label{thm:seedinversion}
On every finite divisor-closed coefficient set, let $A_i$ be the coefficient
of slot $i$ at its anchor cut and let $\kappa_i$ be its canonical seed.  In
vector form,
\begin{equation}
\label{eq:seedmatrix}
 A=(1+K)\kappa,\qquad
 \kappa=(1+K)^{-1}A=\sum_{n\ge0}(-K)^nA.
\end{equation}
In particular,
\begin{equation}
\label{eq:seedchain}
 \kappa_i=A_i+
 \sum_{n\ge1}(-1)^n
 \sum_{i_0<\cdots<i_n=i}
 K_{i_ni_{n-1}}\cdots K_{i_1i_0}A_{i_0},
\end{equation}
where $<$ includes strict divisibility of the deformation monomials and the
intermediate boundary exponents are summed.  Thus every seed is a finite
alternating sum of anchored open FJRW coefficients weighted by ordered
broken-line chains.
\end{theorem}

\begin{proof}
Evaluate Theorem~\ref{thm:transport} at the anchor of $i$.  The seed in slot
$i$ contributes $\kappa_i$, since a positive-degree wall acting on it
overshoots its deformation monomial.  A seed in $j<i$ contributes
$K_{ij}\kappa_j$, and no other seed contributes.  Hence
$A_i=\kappa_i+\sum_{j<i}K_{ij}\kappa_j$.  The kernel strictly increases
deformation monomials and is nilpotent in the finite quotient, so its geometric
series inverse is finite and gives
\eqref{eq:seedmatrix}--\eqref{eq:seedchain}.  \qedhere
\end{proof}

The inversion has a zero-dimensional enumerative lift.  Let
$\mathcal Z_i^{\mathrm{anc}}$ be the signed, automorphism-weighted zero-cycle
of the balanced descendent Witten multisection at $q_i$, so
$\deg\mathcal Z_i^{\mathrm{anc}}=A_i$.  For a term
$c_\lambda u_\lambda X_{k_\lambda}$ in a ray logarithm, let
$\mathcal Z_\lambda^{\mathrm{crit}}$ be the signed weighted union of the zero
loci of the corresponding critical-graph homotopies in the staged
realization.  Exact-support ray purity identifies its degree with $c_\lambda$.

For a decorated bend chain $\gamma$ beginning at $i_0=(E,b)$, define the
rational oriented cascade cycle
\begin{equation}
\label{eq:cascadechain}
 \mathcal C_\gamma=
 \frac{\prod_h\epsilon_{j(h)}
 \det(a^{h-1},k_{\lambda_h}+(1,1))}{\prod_jn_j!}
 \left(\mathcal Z_{i_0}^{\mathrm{anc}}
 \times\prod_h\mathcal Z_{\lambda_h}^{\mathrm{crit}}\right).
\end{equation}
Equations (5.36)--(5.37) of \cite{GKT} identify one such bend with the corresponding
critical boundary zero locus: its zero count is the critical-homotopy count
times the balanced count, multiplied by precisely the determinant/gluing
factor in \eqref{eq:cascadechain}.  Associativity of the assembling operation
\cite[\S2.3.2]{GKT} allows these
factors to be iterated.  Expanding the kernels along a strict divisibility
chain $i_0<\cdots<i_n=i$ therefore gives a finite cascade cycle
$\mathcal C(i_0<\cdots<i_n=i)$ of degree
$K_{i_ni_{n-1}}\cdots K_{i_1i_0}A_{i_0}$.

\begin{theorem}[Corrected relative-Euler cascade class]
\label{thm:seedcascade}
The rational zero-cycle
\begin{equation}
\label{eq:seedcycle}
 \mathcal Z_i^{\mathrm{seed}}
 =\mathcal Z_i^{\mathrm{anc}}
 +\sum_{n\ge1}(-1)^n
   \sum_{i_0<\cdots<i_n=i}\mathcal C(i_0<\cdots<i_n=i)
\end{equation}
is finite and satisfies
\[
 \deg\mathcal Z_i^{\mathrm{seed}}=\kappa_i.
\]
Its class in rational oriented zero-bordism is independent of the endpoint
representatives, canonical homotopies, and transverse perturbations used in
the staged realization.
\end{theorem}

\begin{proof}
Taking degrees in \eqref{eq:seedcycle} gives the right side of
\eqref{eq:seedchain}, hence $\kappa_i$.
Anchored degrees are coefficients of the canonical cut potentials, and the
critical degrees are coefficients of the canonical ray factors; freeness and
ray purity make both independent of the chosen representatives.  Finally,
rational oriented zero-bordism is classified by weighted degree,
$\Omega^{\mathrm{or}}_0\otimes\QQ\cong\QQ$, which proves independence of the
class.  \qedhere
\end{proof}

We next package the separate products in \eqref{eq:seedcycle} into one
geometric object.  A \emph{virtual orbifold with corners} below means a finite
Kuranishi atlas with isotropy: its charts are finite-dimensional
orbifolds-with-corners carrying relatively oriented obstruction bundles,
and a codimension-$c$ coordinate change stabilizes both the base and the
obstruction bundle by the same $c$ normal directions.  A compatible transverse
multisection defines a weighted branched virtual cycle
\cite{MWK}.  All charts used here are products and corners of the
finite-dimensional GKT orbifolds; no additional analytic regularization is
involved.

\begin{lemma}[Polytope-relative canonical extension]
\label{lem:polytopeextension}
Let $P$ be a compact smooth polytope with corners and let $\mathcal G$ be a
finite construction-closed collection of graded graphs.  Suppose compatible
canonical multisection families are prescribed on
\[
 (\partial P\times\mathcal{PM}_\Gamma)
 \cup(P\times\partial\mathcal{PM}_\Gamma)
\]
and near $P\times\partial^+M_\Gamma$, for every $\Gamma\in\mathcal G$.
Assume that the graph-boundary values are assembled from the corresponding
families on the base components and are strongly positive near the positive
boundary.  The families extend, relative to a smaller collar of the
prescribed set, over $P\times\mathcal{PM}_\Gamma$.  They may be chosen
automorphism-invariant and transverse on every parameter face and graph
stratum.  Two extensions fixed on the prescribed set are joined by such an
extension over $P\times[0,1]$.
\end{lemma}

\begin{proof}
Repeat the induction on $\dim\mathcal{PM}_\Gamma$ in
\cite[Lem.~3.39]{GKT}.  Once the lower graph families are fixed, assembling
prescribes the section on
$P\times\partial\mathcal{PM}_\Gamma$; associativity of assembling
\cite[\S2.3.2]{GKT} makes the prescriptions agree at graph corners.
Orbifold collars and the first three steps of the BCT extension argument used
in that proof extend the data while preserving strong positivity.
On the compact complement choose finitely many local multisections spanning
the fibers.  Generic small linear combinations are simultaneously transverse
on the finitely many parameter faces and graph strata.  Taking the finite
automorphism orbit gives an invariant multisection.  The same construction
over $P\times[0,1]$ proves the final assertion. \qedhere
\end{proof}

\begin{lemma}[Iterated critical corner]
\label{lem:multicorner}
Let $\gamma=(\lambda_1,\ldots,\lambda_m)$ be a labeled decorated bend
history acting on an anchored balanced problem.  There is a stable graph
$\Gamma_\gamma$ with $m$ distinguished boundary edges and a distinguished
balanced spine such that $\mathcal{PM}_{\Gamma_\gamma}$ is a
codimension-$m$ corner of the fully smoothed outer moduli space.  Detaching
the distinguished edges gives a finite map
\[
 F_\gamma:\mathcal{PM}_{\Gamma_\gamma}\longrightarrow
 \mathcal{PM}^{\mathrm{bal}}\times
 \prod_{h=1}^m\mathcal{PM}_{\lambda_h}^{\mathrm{crit}}.
\]
The descendent Witten bundle and assembled multisection on the corner are
the pullbacks of the external sums on the right.  These identifications do
not depend on the order of detachment.
\end{lemma}

\begin{proof}
For $m=1$ this is the critical boundary graph of
\cite[Not.~5.6]{GKT}: Definition~2.6 and Equation~(2.7) there give the
codimension-one closed stratum, while Proposition~2.14 and
Observation~3.11 identify the bundle; canonicity gives the section identity
used in (5.37).  Inductively attach the next critical component to the rooted
balanced component.  The labeled realization gives disjoint labels to
repeated insertions, so the graph remains stable.  Proposition~2.4 of
\cite{GKT} supplies the orbifold corner and one smoothing coordinate per
edge.  The assertion follows by iterating detachment; independence of order
is precisely associativity of assembling. \qedhere
\end{proof}

For $j<i$, let $\mathfrak T_{ij}$ be the finite rational virtual
zero-orbifold obtained by taking the union of the decorated bend charts in
the transport from $q_j$ to $q_i$, including the GKT gluing covers, signs,
and symmetric quotients.  Lemma~\ref{lem:multicorner} and
\cite[(5.36)--(5.37)]{GKT} give
\begin{equation}
\label{eq:transportvirtualdegree}
 \deg[\mathfrak T_{ij}]^{\mathrm{vir}}=K_{ij}.
\end{equation}

\begin{theorem}[Compact virtual cascade realization]
\label{thm:virtualcascade}
Fix a finite divisor-closed coefficient set and a target slot $i$.  There is
one compact oriented virtual orbifold with corners
$\overline{\mathfrak C}^{\,\mathrm{bar}}_i$ of virtual dimension zero, one
assembled obstruction atlas on it, and one coherent transverse
multisection $\mathfrak s_i$ such that
\[
 [Z(\mathfrak s_i)]^{\mathrm{vir}}
 =\mathcal Z_i^{\mathrm{seed}}.
\]
In particular its virtual degree is $\kappa_i$.  Its virtual
zero-bordism class is independent of all paths, collars, higher fillers, and
transverse perturbations.  These objects can be chosen compatibly under
enlargement of the coefficient set.
\end{theorem}

\begin{proof}
We first verify the nontrivial composition face.  For
$i_0<i_1<i_2$, use Moore paths in the staged boundary-condition interval.
The concatenated path through $q_{i_1}$ and the direct path from $q_{i_0}$
to $q_{i_2}$ are homotopic relative to their endpoints.  Pull back the GKT
families over this homotopy, use compactified ordered configurations for the
bend times, and apply Lemma~\ref{lem:polytopeextension}.  After the broken
configuration faces are identified by Lemma~\ref{lem:multicorner}, the zero
locus is a compact virtual one-bordism with
\begin{equation}
\label{eq:lengthtwoboundary}
\begin{split}
 \partial[\mathfrak B(i_0,i_1,i_2)]^{\mathrm{vir}}
 ={}&
 [\mathfrak T_{i_2i_1}\times\mathfrak T_{i_1i_0}
   \times\mathcal Z_{i_0}^{\mathrm{anc}}]^{\mathrm{vir}}\\
 &-[\mu_2(\mathfrak T_{i_2i_1},\mathfrak T_{i_1i_0})
   \times\mathcal Z_{i_0}^{\mathrm{anc}}]^{\mathrm{vir}}.
\end{split}
\end{equation}
Here $\mu_2$ concatenates the histories and forgets the waypoint; it retains
the same critical and balanced factors.  Thus
\eqref{eq:lengthtwoboundary} is a cycle-level boundary identity, not only the
numerical equality
$\deg\mu_2=K_{i_2i_1}K_{i_1i_0}$.

The higher coherences are obtained inductively over the multiplihedra
$J_m$, whose faces encode the compositions of consecutive subwords
\cite{MauWoodward}.  The lower families agree on intersecting faces because
Moore-path concatenation and GKT assembling are associative.
Lemma~\ref{lem:polytopeextension} fills
$J_m\times\mathcal{PM}_\Gamma$ relative to its boundary.  Consequently the
finite directed system of transport correspondences and anchored cycles has
a homotopy-coherent relative bar realization.

Its nondegenerate $n$-cells are
\[
 \coprod_{i_0<\cdots<i_n=i}
 \mathfrak T_{i_ni_{n-1}}\times\cdots\times
 \mathfrak T_{i_1i_0}\times\mathcal Z_{i_0}^{\mathrm{anc}}.
\]
Collapse faces compose adjacent transports, breaking faces are the
multi-critical corners of Lemma~\ref{lem:multicorner}, and the remaining
external faces are relative faces with a nonzero normal section.  On an
$n$-cell, cross the product GKT chart with the cellular $n$-handle and add
its $n$ normal directions to the obstruction bundle.  On a face this is the
stabilization by the boundary coordinate.  More explicitly, a breaking-face
coordinate change is the span consisting of the finite labeled detaching
cover and the closed corner embedding of \cite[Def.~2.6]{GKT}; a collapse-face
coordinate change is the collared parameter-face embedding.  Proposition~2.14
and assembling identify the bundle data, and the smoothing coordinate
identifies the tangent and obstruction normal quotients.  Thus the spans
satisfy the Kuranishi index condition and, by the higher face identities,
form a finite atlas with isotropy \cite{MWK}.  The standard cellular Morse
section has one zero of Morse
index $n$ and hence local sign $(-1)^n$.  Keeping the original product
multisections near these zeros and applying
Lemma~\ref{lem:polytopeextension} elsewhere gives
\[
 [Z(\mathfrak s_i)]^{\mathrm{vir}}
 =\mathcal Z_i^{\mathrm{anc}}
 +\sum_{n\ge1}(-1)^n
 \sum_{i_0<\cdots<i_n=i}
 [\mathfrak T_{i_ni_{n-1}}\times\cdots\times
   \mathfrak T_{i_1i_0}\times
   \mathcal Z_{i_0}^{\mathrm{anc}}]^{\mathrm{vir}},
\]
which is \eqref{eq:seedcycle}.

It remains to check the rational weights.  Permuting two virtual
zero-dimensional critical factors changes the base orientation by
$(-1)^{d_1d_2}$ and the obstruction orientation by
$(-1)^{r_1r_2}$; since $d_a=r_a$, their product is $1$.  Thus the
$S_{n_j}$-quotient for $n_j$ same-ray bends is relatively
orientation-preserving and contributes $1/n_j!$.  The GKT boundary
orientation and finite gluing degree give the determinant factor in
\eqref{eq:cascadechain}.  Hence the displayed virtual cycle is the stated
cycle itself.  Its degree is $\kappa_i$ by
Theorem~\ref{thm:seedinversion}.

Strict divisibility bounds the bar length, so the finite realization is
compact.  Applying Lemma~\ref{lem:polytopeextension} after adding one
parameter gives independence by a virtual one-bordism.  Applying its relative
form along the construction-closed exhaustion, with the old charts fixed,
gives coefficientwise compatibility. \qedhere
\end{proof}

\begin{remark}[Virtual versus ordinary presentation]
\label{rem:virtualcascade}
Theorem~\ref{thm:virtualcascade} is the single cascade object naturally
provided by the geometry.  It is stronger than the formal zero-bordism class
of Theorem~\ref{thm:seedcascade}: its cells are attached by the actual GKT
detaching, smoothing, and continuation faces.  It does not assert a further
presentation by one smooth GKT orbifold and one fixed-rank ordinary
orbibundle.  Across an $n$-cell face, one cellular base direction and one
obstruction direction disappear together; the Kuranishi stabilization
records this canonically.  A fixed-rank presentation would require a
separate global thickening theorem and is not needed for the virtual Euler
cycle.
\end{remark}

The mechanism is corroborated computationally. At orders $t^{2}$ and $t^{3}$
the pro-nilpotent wall diagonals are verified in full: their
coefficients, with their cut dependence and signs, are
generated by the axis seeds $x^{r},y^{s}$ bending at the computed wall
functions. In the pure-$t_{3,3}$ family the verification extends through
$t^{4}$, nonlinearly (\S\ref{sec:cert}): the ansatz
$\exp(\tfrac1{50}t_{3,3}^{2}X_{1,1}+c_{4}t_{3,3}^{4}X_{2,2})$ carries
$W^{\nu^{\min}}$ to $W^{\nu^{\max}}$ if and only if $c_{4}=0$, both order-4
coefficients pinning $c_{4}$ independently, and the double bend of the $y^{5}$ seed
computes $\langle\tau_{0}^{(3,3)4}\sigma_{1}^{2}\sigma_{2}^{7}\sigma_{12}
\rangle^{\min}=-\tfrac{12}{25}$ from the wall function alone. The mixed
$t_{2,3}$--$t_{3,3}$ family through $t^{4}$ (the first with noncommuting
walls) passes the same checks (\S\ref{sec:cert}), with order-4
wall coefficients $c_{(1,2)}=-\tfrac1{750}$ (with a contribution from the commutator
$[\theta_{(1,2)},\theta_{(1,1)}]$, the first scattering vertex of the theory)
and $c_{(0,2)}=-\tfrac1{1250}$.

In the off-ray case of Example~\ref{ex:farside} the canonical seed of the
extreme coefficient $(7,1)$ is $-\tfrac{10}{3}$: its normal-form zero is the
superposition of the seed with a \emph{far-side bend} (a surviving foreign
contribution at the coefficient's anchor cut; Example~\ref{ex:farside}),
$0=-\tfrac{10}{3}+\tfrac{10}{3}$ (with the naive anchored-value seed the
transport misses $W$ by the constant $-\tfrac{10}{3}$; with the canonical seed
it agrees exactly in all four chambers of that finite family). On the foreign-free coefficients
(those whose anchor cut receives no foreign bending) the seed coincides with
the anchored value, zero at extremes. Every coefficient on the
$(1,1)$-direction is foreign-free (Theorem~\ref{thm:nobackscatter}), and the
exhaustive enumeration behind Example~\ref{ex:farside} verifies that all
extreme coefficients of the two computed families are foreign-free as well:
those coefficients are purely tropical. This establishes the final clauses of
Theorem~\ref{thm:transport}.  Theorems~\ref{thm:seedcascade}
and~\ref{thm:virtualcascade} supply the zero-dimensional enumerative class of
every seed and its single compact virtual cascade realization.  The
position-space reading of the direction sector remains open.

\begin{observation}
\label{obs:richer}
The open theory is strictly richer than its tropicalization. Broken lines generate
exactly the part of the potential reachable by transport from the seeds; the seed
system is input the
diagram does not determine, and it is not zero. This has no analogue in the Fano
case of \cite{Gross10,CPS}, where every correction is disk-generated. The witness is
Example~\ref{ex:middleseeds}: a middle seed of value $-\tfrac{25}{16}$ that no
bending can produce.
\end{observation}

\begin{example}[The non-tropical part; middle seeds]
\label{ex:middleseeds}
Two things this example shows: bending cannot generate everything, and seeds are
not confined to the coefficients that carry no wall. For the first, the invariant
$\langle\tau_{0}^{(3,3)3}\sigma_{1}^{4}\sigma_{2}^{4}\sigma_{12}\rangle
=\tfrac15$ of \eqref{eq:newinvariants} has no bending source ($t$-monomials
multiply along a broken line, and the unique candidate wall fixes the
$t_{3,3}$-seed exactly), so such coefficients are seeds. For the second, take the $N=2$ singleton tower $t_{3,3,d=2}$ of
$x^{5}+y^{5}$, the exact solution of Theorem~\ref{thm:linearorder} has equal
wall coefficients $c_{1}=c_{2}=\tfrac5{144}$, extreme value
$\langle\tau_{2}^{(3,3)}\sigma_{1}^{3}\sigma_{2}^{13}\sigma_{12}\rangle
=\tfrac{25}{36}$, and \emph{middle seed}
$\langle\tau_{2}^{(3,3)}\sigma_{1}^{8}\sigma_{2}^{8}\sigma_{12}
\rangle^{\mathrm{mid}}=-\tfrac{25}{16}$, the $\Ascr$-condition holding in all
three chambers.
\end{example}

\begin{example}[Far-side bends]
\label{ex:farside}
Theorem~\ref{thm:nobackscatter} excludes foreign bending only from beyond the
target's own direction; arrivals from the $(1,1)$-side survive it. We call a
surviving foreign contribution at a target's anchor cut a \emph{far-side
bend}, after the path shape: the seed lies on the far side of the crossed
wall from the target's cut. Characteristically these are near-axis descendent
walls bending seeds from the linear, flat-coordinate part of the deformation.
Exhaustive enumeration shows the two computed families of
$x^{5}+y^{5}$ admit no such configurations (their verified transport identities
are fully explained by Theorems~\ref{thm:nobackscatter}
and~\ref{thm:linearorder}), but composite diagonals with descendents do. The point of the two smallest cases is that the naive rule fails and the canonical
seed is what repairs it: in the off-ray case the anchored value is $0$ while the
seed is $-\tfrac{10}{3}$, the difference being exactly a far-side bend. In detail,
$J=\{(0,0,2),(1,1,0)\}$ and the off-ray $J=\{(0,0,2),(2,1,0)\}$ are solved
completely (\S\ref{sec:cert}): the
unit-twist tower $t_{0,0,2}$ carries nonzero wall coefficients $\tfrac56$ on both
of its near-axis rays,
the composite diagonals have forced own wall coefficients $1$, and the
$\Ascr$-conditions hold in every chamber of the finite family. In the off-ray case the far-side bend
contributes $\tfrac{10}{3}$ to the \emph{extreme} coefficient $(7,1)$ in its
anchor cut: the anchored value ($0$, the normal form) is not the seed, and
the correct object is canonical.
\end{example}


\subsection{Realization by staged homotopies}
\label{ssec:realization}

\begin{proposition}[Descendent walls accumulate; the primary diagram is locally finite]
\label{prop:localfinite}
Fix $r,s\ge2$ and an interior twist $(a,b)$. Let $d$ range over $\ZZ_{\ge1}$, so that $t_{a,b,d}$ runs over deformation
parameters of total degree one. The rays of the associated critical graphs are
exactly
\[
R(a,b)=\bigl\{(a+jr+1,\ b+\kappa s+1)\ :\ j,\kappa\in\ZZ_{\ge0}\bigr\},
\qquad d=j+\kappa+1,
\]
their wall functions in $\Dscr^{r,s}$ are all nonzero, and their slopes are dense
in $(0,\infty)$. Hence the walls of $\Dscr^{r,s}$
accumulate in every subsector of the direction sector already at total degree one
in the deformation parameters, for every $(r,s)$, the simple singularities
included; those carry no primary wall at all (Theorem~\ref{thm:trichotomy}). By contrast, write
$\Dscr^{r,s}_{0}$ for the \emph{primary diagram}, the walls with
$\mathbf d=\mathbf 0$, and $\mathfrak m$ for the ideal generated by the
deformation parameters. Then $\Dscr^{r,s}_{0}$ has finitely many walls modulo
$\mathfrak m^{n}$ for every $n$, and its chambers are open angular sectors.
\end{proposition}

\begin{proof}
For a singleton $J=\{(a,b)\}$ one has $\ell_{1}=\ell_{2}=0$ and $N=d$, so by
Definition~\ref{def:diagonal} the critical graphs $\Lambda_{J,p}$, $1\le p\le d$,
have boundary degrees $(a+(p-1)r+1,\ b+(d-p)s+1)$; putting $j=p-1$, $\kappa=d-p$
gives $j+\kappa=d-1$ and the stated set, every element of which is attained by
taking $d=j+\kappa+1$. For density, fix $t>0$ and let $j$ be large enough that
$t(a+jr+1)>1+b$; with $\kappa=\lceil (t(a+jr+1)-1-b)/s\rceil$ one has
$\kappa s\in[\,t(a+jr+1)-1-b,\ t(a+jr+1)-1-b+s)$, whence
\[
\Bigl|\frac{b+\kappa s+1}{a+jr+1}-t\Bigr|\ <\ \frac{s}{a+jr+1}\
\xrightarrow[\ j\to\infty\ ]{}\ 0 .
\]
For the wall functions, note that on a singleton diagonal the chamber-index
condition is $\Ascr(J,\mathbf d,\nu)=(-1)^{d}\ne0$, whose $h=1$ part is
$\sum_{p}\omega_{p}\nu_{p}$ with all $\omega_{p}>0$, so $\nu^{\max}_{N}\ne0$;
solving $W^{\nu^{\max}}-W^{\nu^{\min}}=\sum_{p}c_{p}w_{p}$ with
$w_{p}=a_{p}e_{p}-b_{p}e_{p-1}$, $a_{p},b_{p}>0$, gives
$c_{N}=\nu^{\max}_{N}/a_{N}\ne0$ and $c_{p}=c_{p+1}b_{p+1}/a_{p}\ne0$ for every
$p$, and no cancellation between distinct rays is possible at total degree one,
where the monomials are distinct.
For the primary statement, modulo $\mathfrak m^{n}$ only monomials
$u_{J,\mathbf 0}$ with $|J|<n$ occur, and for each such $|J|$ the twists range
over the finite Neveu--Schwarz box, so only finitely many diagonals, hence
finitely many walls, contribute; finitely many rays through the origin cut the
sector into open angular sectors.
\end{proof}

In the statement and proof of the next theorem, call a jump, stage, or degree
\emph{visible} if its class modulo the stated quotient is nonzero; this sense
is unrelated to the NS-visibility of diagonals.

\begin{theorem}[Finite-support and pro-geometric realization;
  $=$ Theorem~\ref{thm:georealization}]
\label{conj:geometric}
Let $S$ be a finite divisor-closed set of deformation monomials containing
$1$, and choose a finite labeled set $I$ with at least $n_{a,b}(S)$ labels of
each twist occurring in $S$, with $n_{a,b}(S)$ as in
Theorem~\ref{thm:georealization}.  The quotient diagram has finitely many
walls, with rays $\rho_1,\dots,\rho_n$ listed in increasing slope.  There is a
canonical homotopy
$(\mathfrak s_{t})_{t\in[0,1]}$, obtained by concatenating families of simple
canonical homotopies in the sense of \cite[\S3.4]{GKT}, from a
family of canonical multisections realizing $\nu_I^{\min}$ to one realizing
$\nu_I^{\max}$, the restrictions of the universal extreme indices, such that
after reduction to the quotient supported on $S$ every nonzero
jump is supported on a single boundary ray, the visible rays occur
in increasing slope as $t$ increases, and the total visible jump on $\rho_{i}$
is the wall automorphism $\theta_{\rho_{i}}$ of
Theorem~\ref{thm:canonical} modulo that quotient.  The choices can be made
compatibly over a cofinal construction-closed exhaustion of all labeled
graph/descendent data, giving the coefficientwise homotopy stated in
Theorem~\ref{thm:georealization}.
\end{theorem}

\begin{proof}
\emph{Step 1: the labeled model, and symmetrization.}
Let $\mathfrak q_S$ be the monomial ideal spanned by monomials outside $S$;
it is an ideal because the complement of a divisor-closed set is upward
closed.  Write a bar for reduction modulo $\mathfrak q_S$. Let
$A_I=\QQ[u_{i,j}:i\in I,\ j\ge0]/(u_{i,j}u_{i,j'})$ be the labeled
coefficient ring of \cite[Def.~4.6]{GKT}, square-zero in each label; write
$u_{A,\mathbf d}=\prod_{i\in A}u_{i,d_i}$ and $u_{i}=u_{i,0}$. The monomials
$u_{J,\mathbf d}$ (indexed by a twist multiset, in the deformation
parameters) and $u_{A,\mathbf d}$ (indexed by a label subset, in $A_{I}$)
live in different rings.
Symmetrization sends, for every descendent level,
\[
t_{a,b,d}\longmapsto
\sum_{\substack{i\in I\\\operatorname{tw}(i)=(a,b)}}u_{i,d}.
\]
Let $\mathfrak q_{I,S}$ be the labeled monomial ideal whose monomials have
symmetric type outside $S$.  Lemma~4.14 of \cite{GKT} says that symmetrization induces
an injection
\[
 \QQ[t_{a,b,d}]/\mathfrak q_S\hookrightarrow A_I/\mathfrak q_{I,S}.
\]
Indeed, a monomial in $S$ maps to the nonempty sum over injective assignments
of its occurrences to labels of the required twists; the saturation bound
$n_{a,b}(S)$ makes the sum nonempty, and distinct symmetric monomials have
disjoint labeled supports. The induced
map on the enlarged Lie algebra preserves shifted boundary rays, including the
marginal $(1,1)$-ray. The action after reduction is still free: the
minimal-support argument of Lemma~\ref{lem:enlarged} detects every nonzero
class supported in $S$ on $W_0$.
Thus the $\mathfrak q_{I,S}$-reductions of $\nu_I^{\min},\nu_I^{\max}$ and of the labeled
diagram identify with their universal symmetric counterparts. Below we
suppress the labeled-lift symbol in the action of a symmetric factor.

\emph{Step 2: the ray factors are symmetric.}
Twist-preserving permutations fix the two symmetric extreme indices, hence by
freeness fix $g^{\min\to\max}$; they also preserve every boundary-ray
subspace. Uniqueness of the factorization in
Theorem~\ref{thm:presentation}(ii), whose filtered proof applies verbatim
over the labeled coefficient ring, therefore makes each
$\theta_{\rho_i}$ symmetric.
Write the boundary-ray factorization, in chronological convention, as
\[
g^{\min\to\max}
=\theta_{\rho_n}\circ\cdots\circ\theta_{\rho_1}
\pmod{\mathfrak q_{I,S}},
\qquad \operatorname{slope}(\rho_1)<\cdots<
\operatorname{slope}(\rho_n).
\]
\emph{Step 3: one stage per ray.}
For each quotient factor choose the symmetric full-group representative
$\widehat\theta_{\rho_i}=\exp(\widehat\eta_{\rho_i})$ whose logarithm is the
finite sum of its terms supported in $S$ and has no terms outside $S$.
Starting with a transverse symmetric canonical family having index
$\nu_I^{\min}$, construct
one stage for each factor. If the current index is $\nu$, then
$\widehat\theta_{\rho_i}\nu$ is again a symmetric element of $\Om(I)$. The
realizability statement $\Om(I)=\Upsilon_{r,s}(I)$
(\cite[Cor.~5.13]{GKT}, together with \cite[Prop.~3.45]{GKT}), read in the
enlarged sense of Lemma~\ref{lem:enlarged} on the marginal ray, supplies a
symmetric canonical endpoint family with this index. The constructive
realization in \cite[Thm.~5.12, Lem.~5.14]{GKT} may be chosen transverse, and
twist symmetrization is performed branchwise, preserving transversality. Join
the two transverse endpoint families by a family of simple canonical homotopies
as in \cite[Lem.~3.39]{GKT}.

\emph{Step 4: ray purity of each stage.}
We claim that every jump internal to this stage is supported on $\rho_i$.
For the actual chronological order of its subset times,
\cite[Thm.~5.8]{GKT} writes the total labeled jump as
\[
\prod_{\varnothing\ne A\subseteq I}^{\mathrm{time}}\exp(v_A),
\]
where $v_A$ has exact label support $A$ (it is the sum over descendent
profiles of terms $u_{A,\mathbf d}X$). In the labeled coefficient ring,
$u_{A,\mathbf d}u_{B,\mathbf e}=0$ when $A\cap B\ne\varnothing$.
For any fixed order this exact-support factorization is unique. Indeed,
a nonzero bracket of disjoint supports has support $A\cup B$; hence
\[
\left.\log\!\left(\prod_A\exp(v_A)\right)\right|_A
=v_A+P_A(v_B:B\subsetneq A),
\]
and induction on $\lvert A\rvert$ recovers every $v_A$.

The total jump and the labeled lift of $\widehat\theta_{\rho_i}$ carry the same initial
endpoint potential to the same final one. Freeness of the enlarged labeled
action (Lemma~\ref{lem:enlarged}) makes them equal. Now induct again on
$\lvert A\rvert$. The exact-$A$ component of the logarithm of the target is
supported on $\rho_i$; by induction all $v_B$ for $B\subsetneq A$ are supported
there as well. Every Lie word in $P_A$ therefore vanishes because the
boundary-ray subalgebra is abelian, as proved in
Theorem~\ref{thm:presentation}(ii). Thus $v_A$ equals the target's exact-$A$
component and is supported on $\rho_i$. This proves the claim, including for
the enlarged torus factor.

\emph{Step 5: the correction stage.}
After the visible stages the resulting full index $\nu^*$ is congruent to
$\nu_I^{\max}$ in the stated quotient. Let $g_{*}$ be the unique full transition from
$\nu^*$ to $\nu_I^{\max}$. Quotient freeness gives $\bar g_{*}=\mathrm{id}$. Join
the current family to a transverse symmetric family realizing $\nu_I^{\max}$ by
one further Lemma-3.39 stage, chosen as above. Its total jump equals $g_{*}$ by full
freeness. Applying the exact-support induction above \emph{after reduction} to
its time-ordered factorization shows successively that every one of its
internal jump factors has trivial class: this correction stage is entirely
invisible modulo $\mathfrak q_{I,S}$.

\emph{Step 6: concatenation.}
Give all stages constant collars at their endpoints, reparametrize them to
successive subintervals of $[0,1]$, and concatenate. Canonicity, strong
positivity and simplicity (\cite[Defs.~3.17, 3.21, 3.35--3.36]{GKT}) survive
this operation: use identical collars and
time reparametrizations simultaneously on every graph and component, which
preserves the boundary-factorization identity \cite[(3.12)]{GKT} and the
component-independence defining simplicity; transversality is unchanged away
from the collars and follows on them from spatial transversality of the
endpoint sections, while strong positivity is slicewise.
The resulting homotopy has the exact endpoints and the asserted visible order
and total jumps. It is not one global Lemma-3.39 normal-form family: the same
marking subset may jump in different ray stages. That distinction is precisely
what evades Proposition~\ref{prop:obstruction}.

\emph{Step 7: compatible choices.}
The correction just used completes a full finite-label family outside $S$ and
therefore should not be nested: a term invisible for $S$ can become visible in
a larger lower set.  Instead exhaust the countable labeled graph/descendent
category by finite \emph{construction-closed} collections
$\mathcal C_1\subset\mathcal C_2\subset\cdots$.  Such a collection contains,
with each datum, all internal-marking subsets, its whole balanced diagonal and
critical predecessors, and every component of $B\Lambda$ required by a
canonical boundary stratum.  Closing finitely many data under these operations
is finite.

Two relative forms of the GKT constructions apply.  First, a transverse
canonical endpoint family on $\mathcal C_n$ extends to
$\mathcal C_{n+1}$, fixed on $\mathcal C_n$, whenever the target chamber index
extends the old one.  To see this, run the proof of
\cite[Thm.~5.12]{GKT} in its order
$(|J|,n(\Gamma),-k_{12}(\Gamma))$ while freezing $\mathcal C_n$.
Lemma~5.14 of \cite{GKT} extends each already prescribed boundary datum, and
the modifications in (5.51)--(5.53) have exact marking set $J$; the observation
following (5.52) says that they do not alter earlier marking sets and that
distinct descendent vectors have independent top-boundary inductions.
Construction-closedness therefore prevents any modification of a frozen
multisection.  Second, a simple canonical homotopy on $\mathcal C_n$ extends
relatively.  The proof of \cite[Lem.~3.39]{GKT} inducts on
$\dim\mathcal{PM}_\Gamma$, prescribes the homotopy on endpoints and canonical
boundary strata by \cite[(3.14)]{GKT}, and perturbs only the remaining
interior.  It consequently leaves every frozen homotopy unchanged.

Choose a cofinal sequence of lower sets $S_n$ detected by the
$\mathcal C_n$.  Fix once and for all pairwise disjoint closed intervals
$J_\rho\subset(0,1)$ ordered by the countable set of wall slopes.  They exist
by inserting the $m$th interval strictly inside the gap left by the finitely
many earlier intervals.  At level $n$, define only the partial endpoint
families on $\mathcal C_n$, realize its finitely many ray factors on the
corresponding $J_\rho$, and leave every other interval constant.  There is no
correction: unseen graph data remain undefined.  Apply the two relative
extensions at level $n+1$; a new ray is constant on $\mathcal C_n$, and a new
term on an old ray extends the old stage because the ray factors are compatible
under restriction.

Organize the exhaustion simultaneously by finite relabeling isomorphism
classes.  When a class first appears, choose the new endpoint data on one
representative, symmetrize it under its finite twist-preserving stabilizer, and
transport it to every later representative; the relative extensions leave all
previously transported data fixed.  This prescription is imposed only on the
endpoints.  The intermediate simple homotopies need not be symmetric because
their distinct subset times generally break that symmetry.

In the union, every graph/descendent datum is eventually fixed literally.
Every such datum depends on only finitely many divisor monomials and boundary
degenerations, hence sees only finitely many nonconstant intervals.  Constant
collars therefore give a genuine smooth canonical homotopy on each
$[0,1]\times\mathcal{PM}_\Gamma$.  All canonical boundary identities hold
because each finite collection of them occurs at a finite stage.  The two
endpoints are natural under every finite twist-preserving relabeling and hence
pro-symmetric.  This is the asserted compatible coefficientwise homotopy; it
does not turn the dense descendent ray set into a locally finite chamber
decomposition.
\end{proof}

A single unstaged homotopy already suffices below the first degree at which one
diagonal carries two walls.

\begin{proposition}[Exact first two-wall order; unstaged realization]
\label{prop:parabolicgeom}
For $m\ge1$, the maximum of the numerical diagonal count
$N(J,\mathbf0)$ over primary multisets with $\lvert J\rvert=m$ is
\[
N_{\max}(r,s;m)
=m+1-\left\lceil\frac{2m}{r}\right\rceil
-\left\lceil\frac{2m}{s}\right\rceil ,
\]
attained by $J=\{(r-2,s-2)^m\}$. Consequently the least primary degree
admitting a diagonal with two walls is
\[
N_2(r,s)=
\min\left\{m\ge3:
\left\lceil\frac{2m}{r}\right\rceil+
\left\lceil\frac{2m}{s}\right\rceil\le m-1\right\},
\]
where the minimum of the empty set is $\infty$.
\begin{enumerate}
\item[(i)] $N_{2}=\infty$ if and only if $\chat_{W}\le1$, that is, if and only if
$W$ is simple or parabolic; when $\chat_W>1$,
$N_{2}\ge\lceil 1/(\chat_{W}-1)\rceil$.
\item[(ii)] When $N_2<\infty$, one \emph{unstaged} Lemma-3.39 normal-form
homotopy realizes the primary diagram modulo $\mathfrak m^{N_{2}}$. When
$N_2=\infty$, such a homotopy exists at every finite order, realizing the
completed primary diagram as the inverse system of these finite-order
statements; no compatibility of the homotopies themselves is asserted.
\end{enumerate}
\end{proposition}

\begin{proof}
Write $\sum_j a_j=r(J)+r\ell_1$ and
$\sum_j b_j=s(J)+s\ell_2$. Since $a_j\le r-2$,
\[
\ell_1\le\left\lfloor\frac{m(r-2)}r\right\rfloor
=m-\left\lceil\frac{2m}{r}\right\rceil,
\]
and similarly for $\ell_2$; both bounds are attained simultaneously by the
maximal twist. Substitution in
$N=\ell_1+\ell_2-m+1$ proves the formula for $N_{\max}$ and hence the displayed
formula for $N_2$. In particular $m=1,2$ never gives $N\ge2$.

The coarser bound from the proof of Theorem~\ref{thm:trichotomy} is
$\ell_{1}+\ell_{2}\le m\chat_{W}$, hence
\begin{equation}
\label{eq:Nbound}
N(J,\mathbf 0)\;\le\;1+m(\chat_{W}-1),
\end{equation}
which is at most $1$ for every $m$ when $\chat_{W}\le1$, and, when
$\chat_W>1$, forces $m\ge1/(\chat_{W}-1)$ whenever $N\ge2$; this is (i)
apart from the converse.
Conversely, if $\chat_{W}>1$, the exact maximal-twist count satisfies
$N_{\max}>m(\chat_{W}-1)-1$, which is at least $2$ for all sufficiently large
$m$, so $N_{2}<\infty$.

For (ii), put $K=N_2$ when $N_2<\infty$ and otherwise fix an arbitrary finite
$K$. Work modulo $\mathfrak q_K=(\mathfrak m^K,\text{descendent parameters})$ and choose
at least $K-1$ labels of every primary twist, so the symmetrization map used in
Theorem~\ref{conj:geometric} is injective in the visible degrees. Then every
visible diagonal has $N\le1$, and a marking subset $A\subseteq I$ determines
the single primary diagonal $(A,\mathbf0)$. Choose transverse symmetric families
realizing $\nu_I^{\min}$ and $\nu_I^{\max}$ and take the one Lemma-3.39 normal-form
family between them. Its jumps occur at the subset times $t_A$. By the element
displayed in the proof of \cite[Thm.~5.8]{GKT}, the reduction of the jump at
$t_A$ is a sum over the critical graphs with marking set $A$: it is zero if
$|A|\ge K$ or no visible wall exists, and otherwise is supported on the unique
wall of that diagonal.

Order the times of all subsets whose visible primary diagonal carries its
unique wall by that wall's slope; place the no-wall and invisible subset times
arbitrarily. A resulting zero jump is harmless. This prescription is legitimate:
\cite[Lem.~3.39]{GKT} constrains the times only to be distinct in $(0,1)$, and its
construction takes them as given: the zeros are moved onto a prescribed
$t_{I(\Gamma)}$ by a diffeomorphism preserving the moduli factor, and the
finitely many non-transverse times are kept off the $t_{A}$ by re-extending
transversally at each prescribed time. Visible jumps sharing a ray commute by the abelianness proved in
Theorem~\ref{thm:presentation}(ii). Full freeness identifies the total product
with $g^{\min\to\max}$; after reduction, its grouping by ray is therefore the
unique boundary-slope factorization
$\prod^{\rightarrow}_{\rho}\theta_{\rho}$. Symmetry and injectivity descend
this labeled statement to the universal primary quotient. Thus the total
visible jump on $\rho_i$ is $\theta_{\rho_i}$.
\end{proof}

\begin{remark}
\label{rem:orderinggap}
The ordering step above is the only place the freedom in the times of
\cite[Lem.~3.39]{GKT} is used, and the freedom is real
rather than assumed: the lemma imposes on the family $(t_{A})_{A\subseteq I}$ only
that its members be distinct points of $(0,1)$, its proof never constructs a $t_{A}$ but consumes
each as a given target, and the induction there runs on $\dim\mathcal{PM}$, which
is unrelated to the marking set indexing the times. What the ordering does
\emph{not} do is separate the walls of a single diagonal: those all carry the one
marking set $J$, so they share the one time $t_{J}$ whatever order is chosen, which
is Proposition~\ref{prop:obstruction}.
\end{remark}

At the first two-wall degree, one global Lemma-3.39 normal-form family provably
does not suffice. This explains why the stagewise construction in
Theorem~\ref{conj:geometric} is necessary.

\begin{proposition}[One normal-form family cannot separate a diagonal's walls]
\label{prop:obstruction}
Let $(J,\mathbf 0)$ be a primary diagonal with $N\ge2$ and let $H^{*}$ be a family
of simple canonical homotopies with times $(t_{A})_{A\subseteq I}$ as in
\cite[Lem.~3.39]{GKT}. Then every possible contribution from the diagonal's
$N$ walls is concentrated at the single time $t_{J}$, with wall $p$ actually
contributing if and only if $\Delta(\Lambda_{J,p})\ne0$. The jump there is
supported on one ray if and only if at most one of these zero counts is nonzero.
In particular no single
such normal-form family gives the ray-by-ray realization of
Theorem~\ref{conj:geometric} once two of the diagonal's walls are actually
crossed.
\end{proposition}

\begin{proof}
The $N$ critical graphs $\Lambda_{J,1},\dots,\Lambda_{J,N}$ all carry the internal
marking set $J$, so by \cite[Lem.~3.39(3)(a)]{GKT} their zeros occur at the one
time $t_{J}$. By the jump formula in the proof of
\cite[Thm.~5.8]{GKT}, the logarithm of the jump there is
\[
\sum_{p=1}^{N}(-1)^{|J|+k_{2}(\Lambda_{J,p})-1}\,\Delta(\Lambda_{J,p})\,
u_{J,\mathbf0}X_{k_{1}(\Lambda_{J,p})-1,\,k_{2}(\Lambda_{J,p})-1}.
\]
Its boundary directions are pairwise distinct: along $p$ the first
boundary degree increases by $r$ and the second decreases by $s$, both by at least
$2$ from values at least $1$, so the slope strictly decreases; and for $N\ge2$ none
of them is the torus direction. The displayed generators lie in distinct
bidegree summands, so no cancellation is possible. Its logarithm is therefore
supported on one ray exactly when at most one coefficient is nonzero.
\end{proof}

The stagewise theorem lies outside this family: it applies Lemma~3.39 separately
to each boundary-ray factor and may therefore reuse the same marking subset at
several stage times. Proposition~\ref{prop:parabolicgeom} succeeds without
staging because when $N\le1$ the at-most-one condition on
$\Delta(\Lambda_{J,p})$ holds for free.

\begin{example}[The first two-ray diagonal of $x^5+y^5$]
\label{ex:tworay}
For $x^{5}+y^{5}$ the order $N_{2}=5$ of
Proposition~\ref{prop:parabolicgeom} is attained by $J=\{(3,3)^{5}\}$,
$\mathbf d=\mathbf 0$: here $m(J,\mathbf 0)=50$, the Euler weight is $25$, and
$N=2$, with balanced profiles $(0,10),(5,5),(10,0)$ and critical graphs of
boundary degrees $(1,6)$ and $(6,1)$: two walls on two \emph{distinct} rays,
the first such diagonal for this Fermat pair. Their jump vectors are
$w_{1}=30e_{1}-5e_{0}$ and $w_{2}=5e_{2}-30e_{1}$. Since $d(J,\mathbf 0)=-3<0$ the
chamber-index condition is $\Ascr(J,\mathbf 0,\nu)=0$, which reads
\[
\tfrac{6}{25}\,\nu_{(0,10)}+\tfrac1{25}\,\nu_{(5,5)}+\tfrac{6}{25}\,\nu_{(10,0)}
\;=\;\tfrac{204}{15625}:
\]
three coefficients against two torsor directions and one condition, as
Theorem~\ref{thm:effective} requires. In the minimal normal form of
Theorem~\ref{thm:canonical} this gives
\[
\bigl\langle(\tau_{0}^{(3,3)})^{5}\,\sigma_{2}^{10}\sigma_{12}\bigr\rangle^{\min}
=\tfrac{34}{625},
\]
with $\bigl\langle(\tau_{0}^{(3,3)})^{5}\sigma_{1}^{10}\sigma_{12}\bigr\rangle^{\max}$
the same by symmetry. (Its agreement with a closed value of
Example~\ref{ex:x5y5} is a numerical accident.) This is the lowest primary degree
for $x^5+y^5$
at which a single unstaged jump would have to carry two rays, hence the smallest
test distinguishing Theorem~\ref{conj:geometric} from one global
Lemma-3.39 normal-form family; every sector computed
in \S\ref{sec:canonical} and in \S\ref{sec:geometric} below has $N\le1$, which is
why the difficulty does not surface there.

The hypothesis of Proposition~\ref{prop:obstruction} is met here. Take $I$ to carry
exactly five markings of twist $(3,3)$, so that this monomial is realized by a single
marking set and the walls above it share one time. Writing
$\nu^{\max}-\nu^{\min}=c_{1}w_{1}+c_{2}w_{2}$ and reading the extreme rows gives
$c_{1}=c_{2}=\tfrac{34}{3125}$, the middle row being consistent; more generally
$c_{1}=\nu_{0}/b_{1}$ and $c_{p+1}=a_{p}c_{p}/b_{p+1}$, so every $c_{p}$ is nonzero as
soon as $\nu_{0}$ is, exactly as in the proof of Proposition~\ref{prop:localfinite}.
No lower-order term can absorb either coefficient: the walls of this tower sit at
$|J|=2,4,5$ and no multiset drawn from the lower orders $|J|=2,4$ sums to $5$, while a wall acting on a lower
coefficient of the potential contributes nothing here, its direction being parallel to
that coefficient's and hence annihilating it (Proposition~\ref{prop:placement}). So both
zero counts are nonzero and both walls are crossed at the one time.

The same computation shows the difficulty is not peculiar to the hyperbolic regime. A
single insertion of maximal twist with two descendents has $N=2$ for every $(r,s)$, with
both coefficients again nonzero, so the obstruction appears in every model once
descendents are admitted, the simple ones included, which carry no primary wall at
all. What is special about $x^{5}+y^{5}$ is that it is the first symmetric
hyperbolic Fermat pair and this is its first primary occurrence.
\end{example}

\begin{corollary}
\label{cor:nosimple}
For $x^{5}+y^{5}$ modulo $\mathfrak m^{6}$, no family of simple canonical homotopies
in one global normal form supplied by \cite[Lem.~3.39]{GKT} between families
realizing the extreme indices $\nu^{\min}$ and $\nu^{\max}$ has all of its jumps
supported on single rays.
\end{corollary}

\begin{proof}
By Example~\ref{ex:tworay} both walls of the diagonal $J=\{(3,3)^{5}\}$ are crossed,
and by Proposition~\ref{prop:obstruction} they share the single time $t_{J}$, so the
jump there meets two distinct rays.
\end{proof}

\begin{remark}
\label{rem:thresholdgeom}
For simple $W$ the \emph{primary} statement carries no walls at all: $N\le0$
throughout by Theorem~\ref{thm:trichotomy}, so $\Dscr^{r,s}_{0}$ is empty and the direction sector is a
single chamber. By Proposition~\ref{prop:localfinite}, though, the descendent
walls of the same model are dense in that same sector. The parabolic pairs are the first
case with primary content, and there the unique
ray is $(1,1)$, carrying the marginal torus direction; the jump it effects is the
one exhibited in \cite[Ex.~4.9]{GKTsurvey}. In the hyperbolic regime diagonals
with $N\ge2$ occur and one unstaged jump can then carry several rays at once.
The exact first degree is given by Proposition~\ref{prop:parabolicgeom}.
Up to the symmetry $(r,s)\leftrightarrow(s,r)$, its complete hyperbolic
classification is
\[
\renewcommand{\arraystretch}{1.15}
\begin{array}{ccc}
\toprule
r&s&N_2(r,s)\\
\midrule
3&7&21\\
3&8&12\\
3&9\le s\le11&9\\
3&s\ge12&6\\
4&5&10\\
4&6\le s\le7&6\\
4&s\ge8&4\\
5&5\le s\le7&5\\
5&s\ge8&4\\
r\ge6&s\ge r&3\\
\bottomrule
\end{array}
\]
Thus $N_2\le21$ for hyperbolic integral Fermat exponents, with the maximum attained only
at $(3,7)$ up to symmetry. The coarser lower bound in
Proposition~\ref{prop:parabolicgeom} is useful for a fixed pair but does not
produce an unbounded sequence approaching the threshold: the exponents are
discrete.
\end{remark}

\begin{remark}[Why staging is necessary]
\label{rem:conjform}
On every finite divisor-closed coefficient set,
Theorem~\ref{conj:geometric} realizes the successive slope cuts by intervals
in $[0,1]$, the two sides of the direction sector by the two endpoints, and the
extreme chamber indices by $\nu^{\min}$ and $\nu^{\max}$.  In a finite primary
quotient these cuts bound actual open chambers; in the descendent completion
the compatible realization is only coefficientwise. The marginal torus factor of
Theorem~\ref{thm:presentation}(ii) is carried by the $(1,1)$-ray, whose fixed
monomials are exactly its own.

Proposition~\ref{prop:obstruction} does not contradict the theorem: it concerns
one global Lemma-3.39 normal-form family, in which a marking subset has only one
jump time. The proof of the theorem instead connects the endpoints of each
boundary-ray factor in a separate stage. Exact labeled-support uniqueness forces
every internal jump of that stage onto the prescribed ray, and concatenation may
reuse a marking subset in later stages. The ray-purity step does not use the
single-critical-graph modification of \cite[Lem.~5.16]{GKT} to prescribe a
stage jump, nor does it identify that lemma's relative Euler number with the
zero count in \cite[Thm.~5.8]{GKT}. (The established endpoint-realizability
theorem \cite[Cor.~5.13]{GKT} remains part of the input.)
\end{remark}

\subsection{Polar extremes and the projective position-space obstruction}
\label{ssec:missing}

The open FJRW invariants of $x^{r}+y^{s}$ are coefficientwise weighted
broken-line transports on the direction sector, with the closed extended invariants
of \S\ref{sec:x5y5} as vertex data.  Theorem~\ref{thm:seedcascade} gives the
seed constants an enumerative zero-cycle class.  We next give the extreme
indices zero-free representatives.  This is an existence statement, not a
uniqueness statement.  We then determine exactly how much position-space
geometry the intrinsic boundary-charge sector of Definition~\ref{def:sector}
supports and prove the obstruction to the stronger interpretation.

Call a canonical family \emph{$\sigma_2$-polar} if, on every wall-bearing
diagonal, its multisection on each non-extreme balanced graph
$\Gamma_{J,p}$, $p\ge1$, is zero-free; call it \emph{$\sigma_1$-polar} if the
same holds for $p\le N(J,\mathbf d)-1$.

\begin{lemma}[Relative nonvanishing]
\label{lem:relative-nonvanishing}
Let $X$ be a connected compact $n$-dimensional orbifold with corners,
$n\ge2$, and let $E\to X$ be a rank-$n$ real orbibundle relatively oriented
with respect to $X$.  Let $A$ be a sufficiently small collar of the full
boundary (and take $A=\varnothing$ when the boundary is empty), and let $s_A$
be a smooth nowhere-zero multisection on $A$.  Then $s_A$
extends, fixed on a smaller collar, to a nowhere-zero smooth multisection of
$E$ if and only if
\[
 \deg e(E;s_A)=0.
\]
The same conclusion holds on a connected compact core of a GKT partial
compactification, with $A$ the collar of its nonexchange and cutoff boundary.
\end{lemma}

\begin{proof}
Choose an orbifold metric, delete the zero section, and normalize radially.
The extension problem is then one for the sphere bundle $S(E)$, with fiber
$S^{n-1}$.  Take a finite relative orbifold triangulation of $(X,A)$
subordinate to the corner strata and bundle charts.  Since $S^{n-1}$ is
$(n-2)$-connected, branches extend through the relative $(n-1)$-skeleton.
On a cell with finite stabilizer, replace a branch by its equally weighted
stabilizer orbit; on overlaps take a common branch refinement.  These
operations preserve nonvanishing.

The only obstruction is consequently the top obstruction
\[
 o(s_A)\in H^n(X,A;\mathfrak o_E\otimes\QQ),
\]
which the usual top-cell clutching calculation identifies with the relative
Euler class.  Poincar\'e--Lefschetz duality, relative orientation, and
connectedness give
\[
 H^n(X,A;\mathfrak o_E\otimes\QQ)
 \cong H_0(X;(\mathfrak o_E\otimes\mathfrak o_X)\otimes\QQ)
 \cong\QQ.
\]
Evaluation on the relative fundamental class is therefore an isomorphism, so
the displayed degree vanishes exactly when the obstruction does.  Extending
over the top cells and smoothing after a common equivariant subdivision gives
the required multisection.  Conversely, a nowhere-zero extension is a
sphere-bundle multisection and has zero relative Euler class.
\end{proof}

\begin{theorem}[Polar representatives of the extremes]
\label{prop:polar}
For every finite internal marking set $I$, the index $\nu_I^{\min}$ admits a
transverse symmetric $\sigma_2$-polar canonical representative and
$\nu_I^{\max}$ admits a transverse symmetric $\sigma_1$-polar canonical
representative.  Conversely, the chamber index of every $\sigma_2$-polar
family is $\nu^{\min}$, and that of every $\sigma_1$-polar family is
$\nu^{\max}$.  The finite choices may be made compatibly over a cofinal
construction-closed exhaustion.
\end{theorem}

\begin{proof}
A zero-free multisection has zero weighted count.  A $\sigma_2$-polar family
therefore has $\nu_{J,p}=0$ for every $p\ge1$ on every wall-bearing diagonal.
The uniqueness part of Theorem~\ref{thm:canonical} identifies it with
$\nu^{\min}$.  The other assertion is symmetric.  Diagonals with $N=0$ impose
no polar condition and are fixed by the chamber-index equations.

It remains to prove existence.  We construct the $\sigma_2$-polar family; the
other construction exchanges the two boundary types and reverses $p$.  Run
the construction of \cite[Thm.~5.12]{GKT} with target $\nu_I^{\min}$, but
refine its extension step as follows.  Fix a wall-bearing diagonal
$(J,\mathbf d)$ and a balanced graph $\Gamma_{J,p}$ with $p\ge1$.  Then
\[
 \nu^{\min}_{J,p}=0,
 \qquad
 n:=\dim\mathcal{PM}_{\Gamma_{J,p}}
 =k_1+k_2+2|J|-2\ge2;
\]
indeed, $J\ne\varnothing$ and $k_1=r(J)+pr\ge r$.

Temporarily label every boundary marking and fix the relative cyclic order of
the $r$-labels and, separately, of the $s$-labels.  The components in this
block are indexed by the binary shuffles of $k_1$ letters $R$ and $k_2$
letters $S$.  Two shuffles differing by one adjacent interchange
$RS\leftrightarrow SR$ meet exchangeable boundary faces.  The graph on
shuffles with these edges is connected.  Glue each exchange pair.  This is
the gluing described in \cite[Rem.~1.54]{GKT}; the exchange map reverses the
induced boundary orientation by \cite[Prop.~2.54]{GKT}.  The result is a
connected relatively oriented orbifold with corners
$\widehat{\mathcal{PM}}_{J,p}$.  Its remaining boundary consists of the
nonexchange boundary and the cutoff near the positive ends.
If $k_2=0$, the shuffle set is a singleton and there are no exchange faces;
the same notation denotes the original connected component and the argument
below applies directly.

The descendent Witten bundle glues as well.  For exchange-equivalent faces
$\Lambda$ and $\Lambda'$, \cite[Obs.~3.9]{GKT} gives
\[
 F_{\Lambda'}\circ\operatorname{Exchange}=F_\Lambda,
\]
so the canonical boundary condition of \cite[Def.~3.21]{GKT} has identical
traces under the gluing.  Before gluing, the zero-cycle on a fully labeled
shuffle block is the signed sum over one lift of every $R/S$ word.  Relabeling
within either boundary type preserves the moduli and bundle orientations
\cite[Lemmas~2.44(2), 2.46(2) and Prop.~2.51(3)]{GKT}.  Consequently, for a
positive rational covering factor $c_{k_1,k_2}$,
\begin{equation}
\label{eq:shuffledegree}
 \deg e(\widehat E;\widehat s)
 =c_{k_1,k_2}\,
 \left\langle
   \prod_{i\in J}\tau_{d_i}^{(a_i,b_i)}
   \sigma_1^{k_1}\sigma_2^{k_2}\sigma_{12}
 \right\rangle^{s,o}
 =0.
\end{equation}
Gluing changes no interior local degree; Proposition~2.54 supplies precisely
the sign needed at the paired faces.

Apply Lemma~\ref{lem:relative-nonvanishing} while holding fixed the already
assembled strongly positive data on the nonexchange boundary and positive-end
collars.  It gives a nowhere-zero multisection on each exchange-completed
shuffle block.  Cutting the block open produces nowhere-zero branches on all
cyclic-order components and matching traces on every exchange pair.  The
three-point exchange bubble is rigid.  Its detaching map identifies an
exchange face, up to its finite label and ghost symmetries, with the smooth
$k_{12}=2$ base.  Take the same weighted finite orbit of the entire collection
of labeled block sections and push it to the unlabeled moduli.  Its seam
restriction is exactly the pullback of the corresponding pushed-down
multisection on the base.  This is a union of branches, not their vector
average, and hence both multisections remain nowhere zero.  Since the
construction was relative at the corners, all higher boundary-assembling
identities are unchanged.

These graphwise choices form one canonical family.  The induction in
\cite[Thm.~5.12]{GKT} orders graphs by the internal marking set, then by
increasing
\[
 n(\Gamma)=s k_1(\Gamma)+r k_2(\Gamma)+(r+s)k_{12}(\Gamma),
\]
and, at fixed $n$, by decreasing $k_{12}$.  When exchange faces are present,
Equation (5.48) there places the $k_{12}=2$ exchange base before its
$k_{12}=1$ balanced graph.  More generally, each additional stable vertex
without internal markings has weighted count at least $r+s$ by the divisibility
relations (5.50) there.  Thus (5.48) shows that a graph depending on this base
has larger $n$, or has the same $n$ and smaller $k_{12}$; it occurs later in
the order.  Treat the base and the exchange completion as one induction block;
Lemma~5.14 of \cite{GKT} supplies
the later extensions.  Changing this exchange datum leaves every chamber
coordinate unchanged: the two contributions have the same pullback by
Observation~3.9 and opposite boundary orientations by Proposition~2.54,
exactly as in \cite[Lem.~5.3]{GKT}.  Thus the ordinary critical-graph
modifications still impose $\nu_I^{\min}$, while the refined extensions make
every $p\ge1$ section zero-free.

Finally take the finite weighted orbit under the twist-preserving
permutations of $I$, as in \cite[Obs.~3.43]{GKT}.  Every orbit branch is
zero-free on the required graphs, so symmetry preserves polarity.  Running
the same relative construction over the construction-closed exhaustion used
in Theorem~\ref{thm:georealization}, freezing all previous collars and
sections, gives the stated coefficientwise compatibility.
\end{proof}

\begin{remark}[Why scalar vanishing is sufficient only after exchange completion]
For fixed traces on every exchange face, the old componentwise warning remains
valid.  For example, on the first two-wall diagonal of $x^5+y^5$ the middle
profile $(k_1,k_2)=(5,5)$ has
\[
 \#\pi_0(\mathcal{PM}_{\Gamma_{J,1}})=\binom{10}{5}=252
\]
by \cite[Rem.~2.27]{GKT}; one zero scalar sum would not force 252 fixed-boundary
relative Euler degrees to vanish.  The theorem does not make that inference.
It varies the exchange traces coherently and joins the 252 shuffle components
before applying the relative obstruction calculation.
\end{remark}

We now turn to the position-space question.  The lattice, cone, axes, ray
placement, determinant pairing and primitive-form shift are intrinsic
boundary-charge data.  The lattice even supplies natural affine
normalizations, for example total boundary charge.  The issue is not the
existence of a coordinate called radius, but whether continuation detects it
and whether the cone supports the required singular affine monodromy.

Put $C_\partial^\times=C_\partial\setminus\{0\}$.  For a finite
divisor-closed coefficient set $S$, write $\Dscr_S$ for the finite union of
wall rays in the quotient.  A \emph{tame path} has off-wall endpoints and is
piecewise smooth and transverse to $\Dscr_S$, apart from finitely many generic
tangencies; tame homotopies are taken relative endpoints and avoid the deleted
vertex.  Let
$\Pi_{\le2}^{\mathrm{tame}}(C_\partial^\times,\Dscr_S)$ denote the resulting
path 2-groupoid.  Let $\mathsf{Can}_S$ be the continuation bicategory whose
objects are transverse canonical multisection families reduced to $S$, whose
one-morphisms are collared concatenations of simple canonical homotopies, and
whose two-morphisms are relative two-parameter canonical families modulo
relative three-parameter families.

\begin{theorem}[Projective coherent sector and no-go theorem;
  $=$ Theorem~\ref{thm:positionobstruction}]
\label{thm:positionobstruction7}
Let $S$ be finite and divisor-closed.
\begin{enumerate}
\item There is a homotopy-coherent pseudofunctor
\[
 \mathcal F_S:
 \Pi_{\le2}^{\mathrm{tame}}(C_\partial^\times,\Dscr_S)
 \longrightarrow\mathsf{Can}_S
\]
which sends a positive elementary crossing of a ray $\rho$ to a geometric
stage with total jump $\theta_{\rho,S}$, and a negative crossing to its reverse.
It is canonical up to coherent equivalence.
\item If
\[
 \pi:C_\partial^\times\longrightarrow\mathbb P(C_\partial)\cong[0,1]
\]
is projectivization, then
\begin{equation}
\label{eq:projectivefactor}
 \mathcal F_S\simeq\overline{\mathcal F}_S\circ\pi_*
\end{equation}
for a coherent transport pseudofunctor on the interval marked by the wall
slopes.  Thus radial paths act by the identity and
$\mathcal F_S\circ(\delta_\lambda)_*\simeq\mathcal F_S$ for every positive
dilation $\delta_\lambda$.  The factorizations are compatible under
restriction and define a coefficientwise pro-functor.
\item The punctured sector has no origin loop.  The total product
$g^{\min\to\max}$ is boundary-to-boundary transport, not holonomy about the
vertex.  The disk fillers in (i) realize backtracking and composition coherence
away from the vertex, but no codimension-two joint comparing two orders of
distinct rays.
\end{enumerate}
Consequently a GKT-determined realization, one using no valuation, period,
stability, or degeneration datum beyond the present continuation theory,
cannot supply a radially detected modulus or a singular affine origin with
monodromy.  In particular the marginal torus factor is not thereby identified
with integral-affine holonomy.  A genuine Gross--Siebert/SYZ position space
requires additional input.
\end{theorem}

\begin{proof}
Write $q=x_1e_1+x_2e_2$ and set
\[
 R(q)=x_1+x_2,
 \qquad
 a(q)=\frac{x_2}{x_1+x_2}.
\]
Then
\begin{equation}
\label{eq:sectorproduct}
 (R,a):C_\partial^\times\xrightarrow{\sim}(0,\infty)\times[0,1],
 \qquad
 (r,t)\longmapsto r((1-t)e_1+te_2).
\end{equation}
If $a_i$ is the slope coordinate of a wall $\rho_i$, then
$\rho_i\setminus\{0\}=(0,\infty)\times\{a_i\}$.  Thus the finite wall
stratification is the product of a radial half-line with a marked interval.
The homotopy
\[
 (r,t)\longmapsto((1-u)r+u,t),\qquad 0\le u\le1,
\]
is a stratum-preserving strong deformation retraction to $R=1$.

The crossing word of a tame path is therefore the crossing word of its angular
projection.  Under a generic homotopy the word changes only by inserting or
deleting $e_i e_i^{-1}$ or $e_i^{-1}e_i$: a tangency to the $i$th marked slope
creates or cancels that pair.  Crossings of distinct slopes cannot exchange
order.  Equivalently, the chamber-adjacency graph is a finite path, hence a
tree, and every word reduces uniquely to the monotone edge path between its
endpoint chambers.

Choose from Theorem~\ref{conj:geometric} the intermediate canonical families
$\mathfrak s_0,\ldots,\mathfrak s_n$ and the ray stages
$h_i:\mathfrak s_{i-1}\to\mathfrak s_i$.  Assign $\mathfrak s_i$ to the
$i$th chamber, $h_i$ to $e_i$, $h_i$ with reversed parameter to $e_i^{-1}$,
and a collared concatenation to a crossing word.  The loop
$h_i*h_i^{-1}$ has compatible canonical data on the boundary of a parameter
disk.  Lemma~\ref{lem:polytopeextension} fills it relative to that boundary.
The same lemma fills relations among several cancellations; the multiplihedra
used in Theorem~\ref{thm:virtualcascade} give all parenthesization coherences.
Applying the relative extension once more joins two choices of representatives,
stages, collars, or fillers.  This proves (i).

The assignment just constructed uses only the angular chamber and the oriented
crossing word.  It is therefore the pullback of the interval construction,
which proves \eqref{eq:projectivefactor}; a radial path has constant projection and
$\pi\delta_\lambda=\pi$.  If $S\subset T$, uniqueness of the ray
factorization gives $\theta_{\rho,T}|_S=\theta_{\rho,S}$, and the
construction-closed relative exhaustion in
Theorem~\ref{conj:geometric} extends the old stages and fillers without changing
them.  Every fixed coefficient sees only finitely many divisor factors.  This
proves (ii), including the pro-statement.

Finally, \eqref{eq:sectorproduct} makes $C_\partial^\times$ contractible.  Its two boundary sides
are not identified, so the slope-ordered path from the $e_1$-side to the
$e_2$-side is not a loop.  Distinct walls are disjoint after deletion of the
vertex, and the crossing-word argument shows that the available disks cannot
compare their two orders.  To create an origin loop one must enlarge the sector
to a surface with an interior puncture or glue its two sides.  Either operation
requires affine transition data, including a linear part in
$\operatorname{GL}(N_\partial)$, which the formal wall automorphisms do not
supply.  This proves (iii) and the final assertion.
\end{proof}

The theorem resolves the intrinsic version of the position-space question
negatively while retaining its strongest positive part: the punctured sector
carries coherent geometric transport, but that transport is projective.  In a
genuine tropical Landau--Ginzburg model, the input includes a singular
integral-affine manifold, a polyhedral decomposition and a multivalued convex
piecewise-linear function, and broken lines use that actual affine geometry
\cite{CPS}.  None of those radial or vertex data is recovered by
\eqref{eq:projectivefactor}.

The same boundary applies to theta functions.  The present results transport
potentials and their seed decomposition, but they do not construct the binary
interacting-broken-line or pair-of-pants operation needed to identify theta
multiplication constants with the open invariants.  Such theta functions (in
the sense of \cite{GHK15,GS-intrinsic}) could still span the Saito--Givental
deformation of $x^r+y^s$ and recast the open mirror theorem
\cite[Thm.~0.3]{GKT} as intrinsic mirror symmetry.  Theorem~\ref{thm:positionobstruction7}
does not rule out that enriched theory; it proves that it requires new input,
for example a Novikov/energy valuation, a period normalization, or a
logarithmic degeneration with specified affine gluing.  The suggested
identification of the marginal torus with monodromy of the elliptic family
underlying Maher's modularity \cite{Maher} is likewise an additional comparison
theorem, not a consequence of the present sector.

\section{Verification}
\label{sec:cert}
All claims are checked by exact polynomial and rational arithmetic (no floating
point); the programs, with replication instructions, are available at
\url{https://github.com/bwuebben/open-fjrw-scattering}. \path{src/gkt_algebra.py}: the algebra of \cite[Def.~4.22]{GKT}
(bracket closure, the classification of Definition~\ref{def:diagonal} against
brute-force enumeration, the first walls of $x^{4}+y^{4}$ and $x^{5}+y^{5}$, the
period-invisibility of the wall action; Theorems~\ref{thm:trichotomy}
and~\ref{thm:effective}(i)). \path{src/diagonal_wall_count.py}: the exact
maximum $N_{\max}$ and the classification of $N_2$ in
Proposition~\ref{prop:parabolicgeom}, checked against direct multiset
enumeration. \path{src/first_two_ray_diagonal.py},
\path{src/two_ray_forced.py}: the diagonal of Example~\ref{ex:tworay}, its
$\Ascr$-condition and normal form, and the nonvanishing of both zero counts
(Corollary~\ref{cor:nosimple}). \path{src/ray_density_exact.py},
\path{src/ray_accumulation.py}: the ray formula and slope density of
Proposition~\ref{prop:localfinite}. \path{src/oscillatory_check.py}: the seed
relations \eqref{eq:seed} from \cite[Thm.~0.3]{GKT} directly
(Remark~\ref{rem:ex55}). \path{src/a_invariants.py}: the labeled-convention
$\Ascr(J,\mathbf d,\nu)$, the convention pin, the tower of
Remark~\ref{rem:tower}. \path{src/canonical_diagram.py}: normal forms and the
wall functions of Theorem~\ref{thm:canonical}
(Example~\ref{ex:walls}), the $t_{3,3}^{4}$ relation, the pure-$\tau$ invariant of
\eqref{eq:newinvariants}. \path{src/scattering.py}: the wall maps
\cite[Lem.~4.20]{GKT} and the $(1,1)$-shift dictionary
(Proposition~\ref{prop:shift}). \path{src/rspin.py}, \path{src/taut_m0n.py},
\path{src/closed_fjrw.py}: the closed data of Example~\ref{ex:x5y5} with its
internal consistency checks. \path{src/verify_thm05.py},
\path{src/verify_thm05_full.sage}: the full numerical verification of
\cite[Thm.~0.5]{GKT}. \path{src/mirror_periods.py}: Lemma~\ref{lem:periods}.
\path{src/anchor_induction.py}, \path{src/backscatter.py}: the randomized
verifications of Theorems~\ref{thm:linearorder} and~\ref{thm:nobackscatter},
the middle-seed tower of Example~\ref{ex:middleseeds}, and the far-side
enumeration. \path{src/transport_engine.py}, \path{src/mixed_sector.py},
\path{src/farside_test.py}, \path{src/canonical_seeds.py}: the family and
far-side verifications of \S\ref{sec:geometric} and the canonical-seed
mechanism of Theorem~\ref{thm:transport}. \path{src/seed_chain_formula.py}:
the determinant weights, exponential factorial, and the $10/3$ far-side term
in Theorems~\ref{thm:seedinversion}--\ref{thm:seedcascade}. All Python
verifications complete
in minutes on commodity hardware; the SageMath verification is the longest
run.

\begin{thebibliography}{GKT26}

\bibitem[BCT24]{BCT21} A.~Buryak, E.~Clader, R.~J.~Tessler,
\emph{Open $r$-spin theory II: The analogue of Witten's conjecture for
$r$-spin disks}, J.\ Differential Geom.\ 128 (2024), no.~1, 1--75;
arXiv:1809.02536.

\bibitem[CPS24]{CPS} M.~Carl, M.~Pumperla, B.~Siebert,
\emph{A tropical view on Landau--Ginzburg models},
Acta Math.\ Sin.\ (Engl.\ Ser.)\ 40 (2024), no.~1, 329--382;
arXiv:2205.07753.

\bibitem[GKT22]{GKT} M.~Gross, T.~L.~Kelly, R.~J.~Tessler,
\emph{Open FJRW theory and mirror symmetry}, arXiv:2203.02435v2 (2022).

\bibitem[GKT26]{GKTsurvey} M.~Gross, T.~L.~Kelly, R.~J.~Tessler,
\emph{Open enumerative geometries for Landau--Ginzburg models},
arXiv:2602.12707v1 (2026).

\bibitem[GHK15]{GHK15} M.~Gross, P.~Hacking, S.~Keel,
\emph{Mirror symmetry for log Calabi--Yau surfaces I},
Publ.\ Math.\ Inst.\ Hautes \'Etudes Sci.\ 122 (2015), 65--168.

\bibitem[GPS10]{GPS10} M.~Gross, R.~Pandharipande, B.~Siebert,
\emph{The tropical vertex}, Duke Math.\ J.\ 153 (2010), 297--362.

\bibitem[Gro10]{Gross10} M.~Gross,
\emph{Mirror symmetry for $\PP^{2}$ and tropical geometry},
Adv.\ Math.\ 224 (2010), 169--245.

\bibitem[GS19]{GS-intrinsic} M.~Gross, B.~Siebert,
\emph{Intrinsic mirror symmetry}, arXiv:1909.07649 (2019).

\bibitem[HLSW22]{HLSW} W.~He, S.~Li, Y.~Shen, R.~Webb,
\emph{Landau--Ginzburg mirror symmetry conjecture},
J.\ Eur.\ Math.\ Soc.\ 24 (2022), no.~8, 2915--2978; arXiv:1503.01757.

\bibitem[KS06]{KS06} M.~Kontsevich, Y.~Soibelman,
\emph{Affine structures and non-Archimedean analytic spaces}, in
\emph{The unity of mathematics}, Progr.\ Math.\ 244, Birkh\"auser (2006), 321--385.

\bibitem[LLSS17]{LLSS} C.~Li, S.~Li, K.~Saito, Y.~Shen,
\emph{Mirror symmetry for exceptional unimodular singularities},
J.\ Eur.\ Math.\ Soc.\ 19 (2017), 1189--1229; arXiv:1405.4530.

\bibitem[Mah24]{Maher} R.~Maher,
\emph{Predictions in open Fan--Jarvis--Ruan--Witten theory via mirror symmetry,
modularity, and wall-crossing}, Ph.D.\ thesis, University of Birmingham, 2024.

\bibitem[MW10]{MauWoodward} S.~Ma'u, C.~Woodward,
\emph{Geometric realizations of the multiplihedron and its complexification},
Compos.\ Math.\ 146 (2010), no.~4, 1002--1028; arXiv:0802.2120.

\bibitem[MW17]{MWK} D.~McDuff, K.~Wehrheim,
\emph{Smooth Kuranishi atlases with isotropy},
Geom.\ Topol.\ 21 (2017), 2725--2809; arXiv:1508.01556.

\bibitem[Wue26]{note} B.~J.~Wuebben,
\emph{The central-charge threshold for wall-crossing in two-variable open FJRW
theory}, companion note,
\texttt{https://github.com/bwuebben/open-fjrw-scattering} (2026).

\end{thebibliography}

\bigskip
\noindent{\sc Bernd J. Wuebben, New York, NY,} \texttt{wuebben@gmail.com}

\end{document}
